\documentclass[preprint,12pt]{elsarticle}

\journal{Journal of Number Theory}

\usepackage[T1]{fontenc}
\usepackage[utf8]{inputenc}
\usepackage{lmodern}
\usepackage{amsmath,amssymb,amsfonts,amsthm,mathtools}
\usepackage{graphicx}
\usepackage{array}
\usepackage[dvipsnames]{xcolor}
\usepackage[pagebackref=true]{hyperref}
\usepackage{bookmark}
\usepackage{pdfcomment}

% Hover previews for internal references (no layout change).
\makeatletter
\let\CG@origref\ref
\let\CG@origeqref\eqref
\newcommand{\CG@gettip}[1]{\@ifundefined{CGTip@#1}{Reference: #1}{\csname CGTip@#1\endcsname}}
\DeclareRobustCommand{\CGref}[1]{\pdftooltip{\CG@origref{#1}}{\CG@gettip{#1}}}
\DeclareRobustCommand{\CGeqref}[1]{\pdftooltip{\CG@origeqref{#1}}{\CG@gettip{#1}}}
\expandafter\gdef\csname CGTip@conj:affine_chowla\endcsname{Conjecture - Standard fixed two-point affine Chowla: Let a\_1,a\_2inmathbb N and b\_1,b\_2inmathbb Z be fixed, with the two affine forms positive for all sufficiently large n and a\_1b\_2-a\_2b\_1ne0. Then sum\_substacknle Y a\_1n+b\_1ge1, a\_2n+b\_2ge1lambda(a\_1n+b\_1)lambda(a\_2n+b\_2)=o(Y).}
\expandafter\gdef\csname CGTip@lem:macroscopic\endcsname{Lemma - Macroscopic subintervals: Assume Conjecture\textasciitilde{}[conj:affine\_chowla] for a fixed pair of affine forms. If 0le A\_Xle B\_X, A\_X/Xto A and B\_X/Xto B with 0le Ale B<infty, then [ sum\_substackA\_X<nle B\_X a\_1n+b\_1ge1, a\_2n+b\_2ge1lambda(a\_1n+b\_1)lambda(a\_2n+b\_2)=o(X). ]}
\expandafter\gdef\csname CGTip@lem:weighted_tao\endcsname{Lemma - Affine logarithmic cancellation with residue-class weight: Let A\_1,A\_2,Qinmathbb N and B\_1,B\_2,Ainmathbb Z be fixed. Assume that the affine forms A\_1k+B\_1 and A\_2k+B\_2 are positive for all sufficiently large k and are non-proportional: A\_1B\_2-A\_2B\_1ne0. If Qk+A>0 for all sufficiently large k, then sum\_substack1le kle K A\_1k+B\_1ge1, A\_2k+B\_2ge1 Qk+Age1 (lambda(A\_1k+B\_1)lambda(A\_2k+B\_2))/(Qk+A) =o(log K).}
\expandafter\gdef\csname CGTip@rem:status_conj\endcsname{Remark - Status of Conjecture\textasciitilde{}[conj:affine\_chowla]: For non-negative intercepts, Conjecture\textasciitilde{}[conj:affine\_chowla] is now [ref], an unrefereed preprint, with a power-of-logarithm saving at every real cutoff; negative intercepts reduce to that case by a fixed shift (Lemma\textasciitilde{}[lem:macro\_rate]). Sections\textasciitilde{}[sec:rational]--[sec:integer] are kept in their original conditional formulation; the unconditional statements are collected in Section\textasciitilde{}[sec:unconditional].}
\expandafter\gdef\csname CGTip@prop:decomposition\endcsname{Proposition - Exact residue-class decomposition: Let (p,q)=1 and Xinmathbb N. Then C\_X (logfrac pq)  = lambda(pq)lfloorfrac Xqrfloor nonumber   + sum\_a=1\textasciicircum{}q-1 sum\_substackkge0 1le qk+ale X pk+lfloor pa/qrfloorge1 lambda(qk+a) lambda (pk+lfloor(pa)/(q)rfloor). For every ane0 the two affine forms in the inner sum are non-proportional.}
\expandafter\gdef\csname CGTip@thm:fixed\endcsname{Theorem - Fixed rational natural resonance: Assume Conjecture\textasciitilde{}[conj:affine\_chowla]. Then frac1X C\_X (logfrac pq)longrightarrowfraclambda(pq)q. For q=1 the exact identity C\_X(log p)=lambda(p)X holds for every Xinmathbb N.}
\expandafter\gdef\csname CGTip@thm:logstep\endcsname{Theorem - Logarithmic microscopic step: Let (p,q)=1 and let cinmathbb R be fixed. Then C\_X\textasciicircum{}log (logfrac pq+frac cX) longrightarrow fraclambda(pq)q mathbf 1\_cge0 unconditionally.}
\expandafter\gdef\csname CGTip@thm:profile\endcsname{Theorem - Rational natural microscopic profile: Assume Conjecture\textasciitilde{}[conj:affine\_chowla].par Let (p,q)=1 and fix cinmathbb R. Then frac1X C\_X (logfrac pq+frac cX) longrightarrow Psi\_p,q(c), where Psi\_p,q(c)= 0, c<0, [1mm] dfraclambda(pq)q, 0le cle q/p, [3mm] (lambda(pq))/(pc), c>q/p. For c>0, Psi\_p,q(c)= lambda(pq)min (frac1q,frac1pc).}
\expandafter\gdef\csname CGTip@cor:integer\endcsname{Corollary - Integer profile and exact right-hand plateau: Assume Conjecture\textasciitilde{}[conj:affine\_chowla]. For fixed mge2, frac1X C\_X (log m+frac cX) longrightarrow 0, c<0, lambda(m), 0le cle1/m, lambda(m)/(mc), c>1/m. Moreover, for every fixed 0<c<1/m, C\_X (log m+frac cX)=lambda(m)X for all sufficiently large X.}
\expandafter\gdef\csname CGTip@thm:naturalmajority\endcsname{Theorem - Unconditional right-hand natural majority window: Let mge1 and 0<c<2/m. Then liminf\_Xtoinfty lambda(m)frac1X C\_X (log m+frac cX) ge 2min (1,frac1mc)-1>0.}
\expandafter\gdef\csname CGTip@thm:logmajority\endcsname{Theorem - Unconditional logarithmic mesoscopic majority: Let (p,q)=1 and let c\_X>0 satisfy (log c\_X)/(log X)longrightarrowtheta, 0letheta<1. Then liminf\_Xtoinfty lambda(pq) C\_X\textasciicircum{}log (logfrac pq+(c\_X)/(X)) ge frac2(1-theta)q-1. In particular, when q=1 and theta<1/2, writing m=p, liminf\_Xtoinfty lambda(m)C\_X\textasciicircum{}log (log m+(c\_X)/(X)) ge1-2theta>0.}
\expandafter\gdef\csname CGTip@eq:o30\endcsname{Equation: C\_X\textasciicircum{}log (log m+(c\_X)/(X))tolambda(m)}
\expandafter\gdef\csname CGTip@eq:o31\endcsname{Equation: Xm(e\textasciicircum{}1/(mX)-1) = 1+frac12mX+O(X\textasciicircum{}-2)>1.}
\expandafter\gdef\csname CGTip@prop:left\endcsname{Proposition - Exact left-hand c-independence interval: Let mge2 and fix -1/m<c<0. For all sufficiently large X, C\_X (log m+frac cX) = sum\_nle Xlambda(n)lambda(mn-1). This states exact independence with respect to c on that finite interval; no asymptotic value of the right-hand sum is asserted.}
\expandafter\gdef\csname CGTip@eq:o33\endcsname{Equation: clongmapsto C\_X (log m+frac cX)}
\expandafter\gdef\csname CGTip@thm:oai\endcsname{Theorem: [Ordinary two-point affine Chowla [ref]] There is an absolute constant kappa>0 with the following property. Let a\_1,a\_2ge1 and b\_1,b\_2ge0 be fixed integers with a\_1b\_2-a\_2b\_1ne0. Then there is C=C\_a\_1,a\_2,b\_1,b\_2 such that |sum\_nle Xlambda(a\_1n+b\_1)lambda(a\_2n+b\_2)|le C fracX(log X)\textasciicircum{}kappa (Xge3, Xinmathbb R).}
\expandafter\gdef\csname CGTip@lem:macro_rate\endcsname{Lemma - Macroscopic blocks with a rate: Let alpha\_1,alpha\_2ge1 and beta\_1,beta\_2inmathbb Z be fixed with alpha\_1beta\_2-alpha\_2beta\_1ne0. Let 0le Ale B<infty be fixed and A\_X=AX+O(1), B\_X=BX+O(1). Then sum\_substackA\_X<kle B\_X alpha\_1k+beta\_1ge1, alpha\_2k+beta\_2ge1lambda(alpha\_1k+beta\_1)lambda(alpha\_2k+beta\_2)=O(fracX(log X)\textasciicircum{}kappa).}
\expandafter\gdef\csname CGTip@thm:profile_uncond\endcsname{Theorem - Rational natural microscopic profile, unconditional: Let (p,q)=1 and fix cinmathbb R. Then frac1X C\_X (logfrac pq+frac cX)=Psi\_p,q(c)+O\_p,q,c((log X)\textasciicircum{}-kappa), with Psi\_p,q as in [eq:profile] and kappa the absolute constant of Theorem\textasciitilde{}[thm:oai].}
\expandafter\gdef\csname CGTip@cor:uncond\endcsname{Corollary - Unconditional forms of Theorem\textasciitilde{}[thm:fixed], Corollary\textasciitilde{}[cor:integer] and Proposition\textasciitilde{}[prop:left]: Let (p,q)=1 and mge2. ; C\_X(log(p/q))=lambda(pq)X/q+O(X(log X)\textasciicircum{}-kappa). ; The limit in Corollary\textasciitilde{}[cor:integer] holds unconditionally, with error term O\_m,c((log X)\textasciicircum{}-kappa). ; For fixed -1/m<c<0 and all sufficiently large X, C\_X (log m+frac cX)=sum\_nle Xlambda(n)lambda(mn-1)=O(frac X(log X)\textasciicircum{}kappa).}
\expandafter\gdef\csname CGTip@rem:trivial_resonance\endcsname{Remark - The trivial resonance is binary Chowla: Take p=q=1 and fix -1<c<0. Since X(1-e\textasciicircum{}c/X)to|c|<1, for Xge X\_0(c) one has 0<n(1-e\textasciicircum{}c/X)<1 for every nle X, hence lfloor e\textasciicircum{}c/Xnrfloor=n-1, and the term n=1 is omitted by definition. Thus, for every Xge X\_0(c), C\_X (frac cX)=sum\_2le nle Xlambda(n-1)lambda(n), while Psi\_1,1(c)=0. The profile of Theorem\textasciitilde{}[thm:profile] at the single point (p,q,c)=(1,1,c) is therefore, for every large X and not only in the limit, the shift-one Chowla sum. No argument that avoids two-point input can produce Theorem\textasciitilde{}[thm:profile]; the residue-class decomposition isolates the main term exactly, and the remainder is a finite set of fixed binary correlations, which Theorem\textasciitilde{}[thm:oai] controls.}
\expandafter\gdef\csname CGTip@lem:floor_k\endcsname{Lemma - Floors on residue classes: For n=Lm+ale X, lfloor e\textasciicircum{}u\_inrfloor=A\_im+b\_a,i+r\_i(n), r\_i(n):=lfloortheta\_a,i+delta\_inrfloor, and |r\_i(n)|le(p\_i)/(q\_i)|c\_i|e\textasciicircum{}|c\_i|+1. In particular r\_i(n) ranges over a finite set R\_i depending only on (p\_i,q\_i,c\_i).}
\expandafter\gdef\csname CGTip@def:classes\endcsname{Definition - Proportionality classes: For fixed (a,r) write f\_a,r,i(m)=mathcal A\_im+mathcal B\_i, let d\_i:=gcd(mathcal A\_i,|mathcal B\_i|) (with gcd(mathcal A,0)=mathcal A), and let (alpha\_i,beta\_i):=(mathcal A\_i,mathcal B\_i)/d\_i, a primitive vector with alpha\_ige1. Declare isim j if (alpha\_i,beta\_i)=(alpha\_j,beta\_j), equivalently mathcal A\_imathcal B\_j=mathcal A\_jmathcal B\_i. For a class mathcal C let g\_mathcal C(m):=alpha\_mathcal Cm+beta\_mathcal C be its primitive form, let Odd(a,r) be the set of classes of odd cardinality and k'(a,r):=|Odd(a,r)|. The block is resonant if k'(a,r)=0. Put varepsilon\_a,r:=lambda(d\_0d\_1cdots d\_k-1)inpm1.}
\expandafter\gdef\csname CGTip@thm:dec_k\endcsname{Theorem - Exact decomposition of the k-point family: For every Xge1 and every choice [eq:uk], C\textasciicircum{}(k)\_X(u)=sum\_substack(a,r) k'(a,r)=0varepsilon\_a,r |I\_a,r(X)| +sum\_substack(a,r) k'(a,r)ge1varepsilon\_a,rsum\_min I\_a,r(X) prod\_mathcal CinOdd(a,r)lambda(g\_mathcal C(m)). The forms g\_mathcal C, mathcal CinOdd(a,r), are primitive, pairwise non-proportional and positive on I\_a,r(X), and k'(a,r)equiv kpmod 2.}
\expandafter\gdef\csname CGTip@lem:len_k\endcsname{Lemma - Block lengths: Let J\_a,r(c)subseteq[0,1] be the set of t with r\_iletheta\_a,i+(p\_ic\_i)/(q\_i)t<r\_i+1 for all i. It is an interval, possibly empty, and |I\_a,r(X)|=frac XL|J\_a,r(c)|+O\_p,q,c(1). More precisely, m: min I\_a,r(X)=(A\_X,B\_X]capmathbb Z with A\_X=AX+O(1), B\_X=BX+O(1) for fixed Ale B.}
\expandafter\gdef\csname CGTip@def:profile_k\endcsname{Definition - k-point profile: Psi\textasciicircum{}(k)\_p,q(c):=frac1Lsum\_substack(a,r) k'(a,r)=0varepsilon\_a,r |J\_a,r(c)| . For odd k there are no resonant blocks, since k'equiv kpmod2, and Psi\textasciicircum{}(k)equiv0.}
\expandafter\gdef\csname CGTip@rem:k2\endcsname{Remark - k=2: For k=2, L=q, and 0sim1 holds if and only if qmid a and r=0, that is (a,r)=(0,0), with d\_0=q, d\_1=p, varepsilon\_0,0=lambda(pq), and |J\_0,0(c)|=min(1,q/(pc)) for cge0, |J\_0,0(c)|=0 for c<0. Definition\textasciitilde{}[def:profile\_k] recovers Psi\_p,q of Theorem\textasciitilde{}[thm:profile], and Theorem\textasciitilde{}[thm:dec\_k] recovers the block structure of its proof.}
\expandafter\gdef\csname CGTip@thm:suf_k\endcsname{Theorem - Sufficiency: In [eq:dec\_k]: ; the pieces with k'(a,r)=1 are O\_p,q,c(Xexp(-c\_0sqrtlog X)) for an absolute c\_0>0; ; the pieces with k'(a,r)=2 are O\_p,q,c(X(log X)\textasciicircum{}-kappa); ; if Ch\_k' holds for every 3le k'le k with k'equiv kpmod2, then frac1XC\textasciicircum{}(k)\_X(u)longrightarrowPsi\textasciicircum{}(k)\_p,q(c) (Xtoinfty).}
\expandafter\gdef\csname CGTip@prop:criterion\endcsname{Proposition - Exact criterion for unconditional closure: Fix (p,q,c). If every block (a,r) with |J\_a,r(c)|>0 has k'(a,r)le2, then [eq:limit\_k] holds unconditionally, with error O((log X)\textasciicircum{}-kappa). Conversely, if some block with |J\_a,r(c)|>0 has k'(a,r)ge3, then Theorem\textasciitilde{}[thm:suf\_k](1),(2) do not cover it, and [eq:limit\_k] at (p,q,c) is equivalent to the statement that the sum over all such blocks is o(X); it does not follow that each block separately satisfies its k'-point estimate.}
\expandafter\gdef\csname CGTip@ex:k4\endcsname{Example - k=4, closed unconditionally: Take (p\_i/q\_i)=(2,3,6) and c=(1,c\_2,1) with 0<c\_2<tfrac13. Then L=1, a=0, r\_2=0 on [0,1], r\_1=lfloor2trfloor and r\_3=lfloor6trfloor. The blocks in t=n/X, with the proportionality classes of the forms (m, 2m+r\_1, 3m, 6m+r\_3) indexed 0,1,2,3, are: small hline t   r   classes   piece hline [0,tfrac16)   (0,0,0)   0,1,2,3   resonant, varepsilon=lambda(36)=1 [tfrac16,tfrac13)   (0,0,1)   0,1,2,3   lambda(6)lambda(m)lambda(6m+1), k'=2 [tfrac13,tfrac12)   (0,0,2)   0,1,2,3   lambda(12)lambda(m)lambda(3m+1), k'=2 [tfrac12,tfrac23)   (1,0,3)   0,2,1,3   resonant, varepsilon=lambda(9)=1 [tfrac23,tfrac56)   (1,0,4)   0,2,1,3   lambda(6)lambda(2m+1)lambda(3m+2), k'=2 [tfrac56,1)   (1,0,5)   0,2,1,3   lambda(3)lambda(2m+1)lambda(6m+5), k'=2 hline By Proposition\textasciitilde{}[prop:criterion], unconditionally, frac1X C\textasciicircum{}(4)\_X(log2+frac1X, log3+fracc\_2X, log6+frac1X)=frac13+O((log X)\textasciicircum{}-kappa). The second plateau ...}
\expandafter\gdef\csname CGTip@thm:integer_k\endcsname{Theorem - Integer resonances at level k: Let m\_1,dots,m\_k-1ge1 be integers and u\_i=log m\_i+c\_i/X. ; If 0<c\_i<1/m\_i for all i, then for Xge X\_0(m,c) C\textasciicircum{}(k)\_X(u)=lambda(m\_1cdots m\_k-1)sum\_nle Xlambda(n)\textasciicircum{}k= lambda(m\_1cdots m\_k-1)lfloor Xrfloor,  k even, [1mm] O(Xe\textasciicircum{}-c\_0sqrtlog X),  k odd. ; If k is even and c\_i>0 for all i, then liminf\_Xtoinfty lambda(m\_1cdots m\_k-1)fracC\textasciicircum{}(k)\_X(u)X ge 2min(1,frac1max\_im\_ic\_i)-1 . ; If -1/m\_i<c\_i<0 for all i, then for Xge X\_0(m,c) C\textasciicircum{}(k)\_X(u)=sum\_substacknle X m\_in-1ge1 forall ilambda(n)prod\_i=1\textasciicircum{}k-1lambda(m\_in-1). No Chowla input is used.}
\expandafter\gdef\csname CGTip@cor:level3_23\endcsname{Corollary - Level three at (2,3) is ordinary three-point Chowla: For every c=(c\_1,c\_2)in(-tfrac12,0)times(-tfrac13,0) and every Xge X\_0(c), C\textasciicircum{}(3)\_X(log2+fracc\_1X, log3+fracc\_2X) =sum\_nle Xlambda(n)lambda(2n-1)lambda(3n-1)  =sum\_substack3le Mle 6X-3 Mequiv3 (mod 6)lambda(M)lambda(M+1)lambda(M+3), while Psi\textasciicircum{}(3)equiv0. Consequently the level-three profile on this open set of c, for all large X, is equivalent to sum\_substackMle Y Mequiv3 (mod 6)lambda(M)lambda(M+1)lambda(M+3)=o(Y).}
\expandafter\gdef\csname CGTip@thm:nec_k\endcsname{Theorem - Necessity: Fix kge2 and a\_1,dots,a\_k-1ge1. For rinmathbb Z\_ge0\textasciicircum{}k-1 put P\_r(Y):=sum\_mle Ylambda(m)prod\_i=1\textasciicircum{}k-1lambda(a\_im+r\_i), pi\_r:=varepsilon\_0,r mathbf 1[k'(0,r)=0], with the classes of Definition\textasciitilde{}[def:classes] for (p\_i,q\_i)=(a\_i,1). Suppose that [eq:limit\_k] holds for (p\_i,q\_i)=(a\_i,1) and every cin(0,infty)\textasciicircum{}k-1. Then P\_r(Y)=pi\_rY+o(Y) for every rinmathbb Z\_ge0\textasciicircum{}k-1. In particular, if the forms m,a\_1m+r\_1,dots,a\_k-1m+r\_k-1 are pairwise non-proportional, then P\_r(Y)=o(Y).}
\expandafter\gdef\csname CGTip@thm:log_k\endcsname{Theorem - Odd levels, logarithmic average: Let kge3 be odd. For every rational resonance [eq:uk] and every cinmathbb R\textasciicircum{}k-1, C\textasciicircum{}(k),log\_X(u)to0.}
\expandafter\gdef\csname CGTip@prop:barrier_poly\endcsname{Proposition - Polynomial-phase barrier, all kge3: Let kge3, 0=h\_0<h\_1<dots<h\_k-1, V:=prod\_0le j<lle k-1(h\_l-h\_j) and R\_i:=Vprod\_jne i(h\_i-h\_j)\textasciicircum{}-1inmathbb Zsetminus0. Let beta be irrational and b\_i(m):=e(beta R\_im\textasciicircum{}k-1). Then: ; for every integer Dge1, uniformly in alphainmathbb R, limsup\_Ytoinftyfrac1DYint\_0\textasciicircum{}Y|sum\_y<mle y+Db\_i(m)e(alpha m)| dyle D\textasciicircum{}-1/2; ; prod\_i=0\textasciicircum{}k-1b\_i(z+h\_iw)=e(beta Vw\textasciicircum{}k-1) for all z,w; hence with c\_w:=e(-beta Vw\textasciicircum{}k-1)mathbf 1[w is P\_0-rough], frac1Ysum\_zle Yb\_0(z)sum\_M<wletau M(c\_w)/(w)prod\_i=1\textasciicircum{}k-1b\_i(z+h\_iw)=fraclfloor YrfloorYsum\_substackM<wletau M w P\_0-roughfrac1w . For P\_0le Mle P\_0\textasciicircum{}2/tau, a range contained in that of [eq:oai124], every P\_0-rough integer of (M,tau M] is prime, and the right-hand side of [eq:barrier\_poly\_E] is sum\_M<pletau M1/pge(tau-1)/(3taulog(tau M))gg\_tau B\textasciicircum{}-1+varepsilon for large B, which exceeds B\textasciicircum{}-1-varepsilon by the factor B\textasciicircum{}2varepsilon. Thus [eq:E\_k] does not follow from [eq:oai127], the ...}
\expandafter\gdef\csname CGTip@prop:barrier_pos\endcsname{Proposition - Barrier for non-negative coefficients, k=3: Let 0<h\_1<h\_2, rho:=h\_1h\_2(h\_1-h\_2), tauin(1,2), H,D\textasciicircum{}*ge1, and let q be a prime with qgemax(2h\_2+1, 7h\_2(D\textasciicircum{}*+1)\textasciicircum{}3, 6|rho|tau\textasciicircum{}2H\textasciicircum{}2). Put b\_0(m)=e((h\_1-h\_2)m\textasciicircum{}2/q), b\_1(m)=e(h\_2m\textasciicircum{}2/q), b\_2(m)=e(-h\_1m\textasciicircum{}2/q). Then: ; for integers 1le Dle D\textasciicircum{}* and every alpha, limsup\_Ytoinftyfrac1DYint\_0\textasciicircum{}Y|sum\_y<mle y+Db\_i(m)e(alpha m)|dyle(3+log q)/(D); ; for Mle H and every c\_wge0 supported on (M,tau M], liminf\_Ytoinfty Re frac1Ysum\_zle Yb\_0(z) sum\_M<wletau M(c\_w)/(w)b\_1(z+h\_1w)b\_2(z+h\_2w)  ge frac12sum\_M<wletau M(c\_w)/(w).}
\expandafter\gdef\csname CGTip@def:local_gowers\endcsname{Definition - Local Gowers norm: For an integer Dge2 let N\_D be the least prime >4D. For b:mathbb Ntomathbb C and yge0 define G\_y:mathbb Z\_N\_Dtomathbb C by G\_y(x)=b(lfloor yrfloor+x) for xin1,dots,2D+1 and G\_y(x)=0 for the other residues, and put |b|\_U\textasciicircum{}s(y;D):=|G\_y|\_U\textasciicircum{}s(mathbb Z\_N\_D).}
\expandafter\gdef\csname CGTip@lem:vonneumann\endcsname{Lemma - Generalised von Neumann inequality: Let kge2, 0=h\_0<h\_1<dots<h\_k-1, N>h\_k-1 prime, and g\_0,dots,g\_k-1:mathbb Z\_Ntomathbb C with |g\_i|le1. Then |mathbb E\_z,winmathbb Z\_Nprod\_i=0\textasciicircum{}k-1g\_i(z+h\_iw)|le|g\_k-1|\_U\textasciicircum{}k-1.}
\expandafter\gdef\csname CGTip@prop:local_control\endcsname{Proposition - Local U\textasciicircum{}k control of the k-point centering error: Let kge3, 0=h\_0<h\_1<dots<h\_k-1, tauin(1,2), Mge1, mathcal J=(M,tau M], |c\_w|le1 for winmathcal J, varpi:=sum\_winmathcal J|c\_w|/w, and D:=lceil h\_k-1tau Mrceil. There is C=C(h\_k-1) such that for all b\_0,dots,b\_k-1:mathbb Ntomathbb C with |b\_i|le1 and all Yge1,  frac1Y|sum\_zle Yb\_0(z)sum\_winmathcal Jfracc\_wwprod\_i=1\textasciicircum{}k-1b\_i(z+h\_iw)|   le C (varpilog(3M))\textasciicircum{}1/2 frac1Yint\_0\textasciicircum{}Y|b\_k-1|\_U\textasciicircum{}k(y;D) dy   +((2D+1)varpi)/(Y). No hypothesis on the support of c\_w or on b\_0,dots,b\_k-2 is needed.}
\expandafter\gdef\csname CGTip@hyp:LU\endcsname{Hypothesis - LU\_k(varrho,K): quantitative local Gowers uniformity of lambda: Let kge3. For every integer Dge2, every finite set of primes mathcal P and every multiplicative b:mathbb Ntomathbb C with |b|le1 and b(p\textasciicircum{}j)=lambda(p\textasciicircum{}j) for all pnotinmathcal P, jge1, limsup\_Ytoinftyfrac1Yint\_0\textasciicircum{}Y|b|\_U\textasciicircum{}k(y;D) dyle K(2+sum\_pinmathcal Pfrac1p)\textasciicircum{}K(log D)\textasciicircum{}-A.}
\expandafter\gdef\csname CGTip@thm:centering_k\endcsname{Theorem - k-point centering under LU\_k: Let kge3 and assume LU\_k(varrho,K) for some K and some varrho>1+3varepsilon, with varepsilon=10\textasciicircum{}-4 as in [eq:oai\_params]. Then [eq:E\_k] holds for the twists b\_i of [ref], any 1-bounded b\_0,dots,b\_k-2, and all coefficients |c\_w|le1 supported on P\_0-rough integers of (M,tau M] with P\_0/taule Mle e\textasciicircum{}C\_0B, uniformly, for all sufficiently large B.}
\expandafter\gdef\csname CGTip@hyp:G\endcsname{Hypothesis - G\_k(h): k-linear divisor graph estimate: Fix 0=h\_0<h\_1<dots<h\_k-1, 1<tau<2, C\_0ge1, T>0. For every sufficiently large B, every admissible family mathcal D, all cutoffs w\_d as in [ref], all |a\_d|le1 and all F\_0,dots,F\_k-1:mathbb Ntomathbb D, limsup\_Xtoinftyfrac1X|sum\_xle Xsum\_dinmathcal Da\_dprod\_i=0\textasciicircum{}k-1F\_i(x+h\_id) g\textasciicircum{}(k)\_d(x)|ll L\_0B\textasciicircum{}-1-c\_*/2, with implied constant depending only on h,tau,C\_0,T and [eq:oai\_params].}
\expandafter\gdef\csname CGTip@thm:conditional_k\endcsname{Theorem - Conditional k-point Chowla in progressions: Let kge3 and 0=h\_0<dots<h\_k-1. Assume LU\_k(varrho,K) with varrho>1+3varepsilon and G\_k(h). Then for all fixed lge1 and binmathbb Z, frac1Nsum\_substackmle N mequiv b (mod l)prod\_i=0\textasciicircum{}k-1lambda(m+h\_i)longrightarrow0 (Ntoinfty).}
\expandafter\gdef\csname CGTip@cor:family_conditional\endcsname{Corollary - The k-point family under LU and G: Let K\_0ge3. Assume, for every 3le kle K\_0, LU\_k(varrho,K) with varrho>1+3varepsilon and G\_k(h) for every h. Then Ch\_k' holds for 3le k'le K\_0, and [eq:limit\_k] holds for every kle K\_0, every rational resonance and every c. In particular [eq:chowla013] holds.}
\expandafter\gdef\csname CGTip@rem:channels\endcsname{Remark - Channel form at a prime: Fix a prime p and shifts H=h\_0,dots,h\_k-1, and write C\_H(X)=frac1Xsum\_nle Xprod\_rlambda(n+h\_r). Complete multiplicativity gives lambda(n+h\_r)=(-1)\textasciicircum{}v\_p(n+h\_r)lambda(U\_p,r(n)), where U\_p,r(n) is the p-free part of n+h\_r, hence prod\_rlambda(n+h\_r)=eta\_p,H(n)b\_p,H(n) with eta\_p,H(n)=(-1)\textasciicircum{}sum\_rv\_p(n+h\_r) and b\_p,H(n)=prod\_rlambda(U\_p,r(n)). The finite-conductor encoding of [ref] (compendium entry 48.0.35) states that for every Jge1 C\_H(X)=mu\_p(H)T\_p,H,0(X)+sum\_j=1\textasciicircum{}Jsum\_substackubmod p\textasciicircum{}j pnmid uwidehatDelta\_p,H,j(u) T\_p,H,j,u(X)+O(kp\textasciicircum{}-J+frac kX), with channels T\_p,H,j,u(X)=frac1Xsum\_nle Xb\_p,H(n)e(un/p\textasciicircum{}j), and that C\_H(X)to0 as soon as every channel tends to 0 (entry 48.0.36). Since b\_p,H=eta\_p,Hprod\_rlambda(cdot+h\_r) and eta\_p,H coincides with a function of period p\textasciicircum{}J' outside a set of density O(kp\textasciicircum{}-J'), each channel is, up to O(kp\textasciicircum{}-J'), a finite combination of the progression sums ...}
\expandafter\gdef\csname CGTip@rem:status\endcsname{Remark - Status of the two hypotheses: Proposition\textasciitilde{}[prop:local\_control], Lemma\textasciitilde{}[lem:vonneumann] and Propositions\textasciitilde{}[prop:barrier\_poly]--[prop:barrier\_pos] are unconditional; Theorem\textasciitilde{}[thm:centering\_k] is the implication LU\_kRightarrow[eq:E\_k], and Theorem\textasciitilde{}[thm:conditional\_k] is conditional on both hypotheses. At fixed D, LU\_k is a statement about averaged 2\textasciicircum{}k-point correlations of b at shifts le3D (expand |G\_y|\_U\textasciicircum{}k\textasciicircum{}2\textasciicircum{}k); it is therefore not a statement obviously weaker than k-point Chowla, and Theorem\textasciitilde{}[thm:conditional\_k] should be read as isolating exactly which uniformity the bias method needs. LU\_k concerns windows of fixed length D while Ytoinfty, with an ordinary average in y and the rate (log D)\textasciicircum{}-A. Matomäki--Radziwill--Tao proved local Fourier uniformity, with the supremum over frequencies inside the integral, for windows Hge X\textasciicircum{}theta with any fixed theta>0 [ref]; Tao showed that the logarithmically averaged Chowla ...}
\expandafter\gdef\csname CGTip@eq:o34\endcsname{Equation: p/q   C\_X/X   lambda(pq)/q   (C\_X-Xlambda(pq)/q)/sqrt X hline 3/2   0.4999230   0.5000000   -0.109 5/2   0.4995180   0.5000000   -0.682 7/2   0.4996260   0.5000000   -0.529 2/3   0.3326445   0.3333333   -0.974 5/3   0.3324220   0.3333333   -1.289 4/3   -0.3322400   -0.3333333   1.546 5/4   -0.2492080   -0.2500000   1.120}
\expandafter\gdef\csname CGTip@eq:o35\endcsname{Equation: c   C\_X/X   Psi\_3,2(c) hline -0.5   -0.000265   0 -0.1   -0.000321   0 0.0   0.499923   0.500000 0.3   0.499923   0.500000 0.6   0.499904   0.500000 0.8   0.415884   0.416667 1.2   0.276538   0.277778 2.0   0.166553   0.166667}
\expandafter\gdef\csname CGTip@eq:o36\endcsname{Equation: Xm(e\textasciicircum{}c/X-1)=1.000000250000041666ldots, C\_X=-999998.}
\expandafter\gdef\csname CGTip@eq:o38\endcsname{Equation: theta   c\_X   lambda(2)C\_X\textasciicircum{}log   2B\_X-1   1-2theta hline 0.00   1.0000   0.958453098   0.916965031   1.000000000 0.20   25.1189   0.765504459   0.530794400   0.600000000 0.40   630.9573   0.573145613   0.144624375   0.200000000 0.45   1412.5375   0.523535345   0.048073466   0.100000000}
\expandafter\gdef\csname CGTip@eq:o1\endcsname{Equation: C\_X(u):=sum\_substack1le nle X lfloor e\textasciicircum{}u nrfloorge1 lambda(n)lambda(lfloor e\textasciicircum{}u nrfloor).}
\expandafter\gdef\csname CGTip@eq:o2\endcsname{Equation: (1)/(X)C\_X (logfrac pq+frac cX) longrightarrow 0, c<0, lambda(pq)/q, 0le cle q/p, lambda(pq)/(pc), c>q/p.}
\expandafter\gdef\csname CGTip@eq:o3\endcsname{Equation: ell=frac qplfloor(pell)/(q)rfloor.}
\expandafter\gdef\csname CGTip@eq:o4\endcsname{Equation: u\_X=log(p/q)+c/X,}
\expandafter\gdef\csname CGTip@eq:o5\endcsname{Equation: sum\_substacknle Y a\_1n+b\_1ge1, a\_2n+b\_2ge1lambda(a\_1n+b\_1)lambda(a\_2n+b\_2)=o(Y).}
\expandafter\gdef\csname CGTip@eq:CXdef\endcsname{Equation: C\_X(u):= sum\_substack1le nle X lfloor e\textasciicircum{}u nrfloorge1 lambda(n)lambda(lfloor e\textasciicircum{}u nrfloor).}
\expandafter\gdef\csname CGTip@eq:o6\endcsname{Equation: H\_X:=sum\_nle Xfrac1n}
\expandafter\gdef\csname CGTip@eq:logdef\endcsname{Equation: C\_X\textasciicircum{}log(u):= frac1H\_X sum\_substack1le nle X lfloor e\textasciicircum{}u nrfloorge1 fraclambda(n)lambda(lfloor e\textasciicircum{}u nrfloor)n.}
\expandafter\gdef\csname CGTip@eq:weighted_tao_nonprop\endcsname{Equation: A\_1B\_2-A\_2B\_1ne0.}
\expandafter\gdef\csname CGTip@eq:weighted_tao\endcsname{Equation: sum\_substack1le kle K A\_1k+B\_1ge1, A\_2k+B\_2ge1 Qk+Age1 (lambda(A\_1k+B\_1)lambda(A\_2k+B\_2))/(Qk+A) =o(log K).}
\expandafter\gdef\csname CGTip@eq:tao_theorem_form\endcsname{Equation: sum\_K/omega(K)<kle K (lambda(a\_1k+b\_1)lambda(a\_2k+b\_2))/(k) =o(logomega(K)).}
\expandafter\gdef\csname CGTip@eq:tao_weight_one_over_k\endcsname{Equation: sum\_kle K (lambda(A\_1k+B\_1)lambda(A\_2k+B\_2))/(k) =o(log K).}
\expandafter\gdef\csname CGTip@eq:residue_weight_expansion\endcsname{Equation: frac1Qk+A=frac1Q frac1k+O\_Q,A(k\textasciicircum{}-2).}
\expandafter\gdef\csname CGTip@eq:decomp\endcsname{Equation: C\_X (logfrac pq)  = lambda(pq)lfloorfrac Xqrfloor nonumber   + sum\_a=1\textasciicircum{}q-1 sum\_substackkge0 1le qk+ale X pk+lfloor pa/qrfloorge1 lambda(qk+a) lambda (pk+lfloor(pa)/(q)rfloor).}
\expandafter\gdef\csname CGTip@eq:o7\endcsname{Equation: lfloorfrac pq(qk+a)rfloor = pk+lfloor(pa)/(q)rfloor.}
\expandafter\gdef\csname CGTip@eq:o8\endcsname{Equation: frac1X C\_X (logfrac pq)longrightarrowfraclambda(pq)q.}
\expandafter\gdef\csname CGTip@eq:o9\endcsname{Equation: C\_X(log p)=lambda(p)X}
\expandafter\gdef\csname CGTip@eq:logstep\endcsname{Equation: C\_X\textasciicircum{}log (logfrac pq+frac cX) longrightarrow fraclambda(pq)q mathbf 1\_cge0}
\expandafter\gdef\csname CGTip@eq:logstep_Y\endcsname{Equation: Y=X\textasciicircum{}1-varepsilon.}
\expandafter\gdef\csname CGTip@eq:logstep_small_perturbation\endcsname{Equation: frac pq(e\textasciicircum{}c/X-1)n=O\_p,q,c(X\textasciicircum{}-varepsilon).}
\expandafter\gdef\csname CGTip@eq:logstep_fractional_separation\endcsname{Equation: (pa)/(q)in [frac1q,1-frac1q].}
\expandafter\gdef\csname CGTip@eq:logstep_prop_initial\endcsname{Equation: (lambda(pq))/(H\_X) sum\_kle Y/qfrac1qk = (H\_Y)/(H\_X)fraclambda(pq)q+o(1).}
\expandafter\gdef\csname CGTip@eq:logstep_tail_mass\endcsname{Equation: (H\_X-H\_Y)/(H\_X)=varepsilon+o(1).}
\expandafter\gdef\csname CGTip@eq:logstep_bilateral\endcsname{Equation: | C\_X\textasciicircum{}log (logfrac pq+frac cX) -(H\_Y)/(H\_X)fraclambda(pq)q | le (H\_X-H\_Y)/(H\_X)+o(1).}
\expandafter\gdef\csname CGTip@eq:logstep_negative_zero_class\endcsname{Equation: lfloor pe\textasciicircum{}c/Xkrfloor=pk-1}
\expandafter\gdef\csname CGTip@eq:logstep_reindex\endcsname{Equation: qk=qell+q, pk-1=pell+(p-1).}
\expandafter\gdef\csname CGTip@eq:profilelimit\endcsname{Equation: frac1X C\_X (logfrac pq+frac cX) longrightarrow Psi\_p,q(c),}
\expandafter\gdef\csname CGTip@eq:profile\endcsname{Equation: Psi\_p,q(c)= 0, c<0, [1mm] dfraclambda(pq)q, 0le cle q/p, [3mm] (lambda(pq))/(pc), c>q/p.}
\expandafter\gdef\csname CGTip@eq:o10\endcsname{Equation: Psi\_p,q(c)= lambda(pq)min (frac1q,frac1pc).}
\expandafter\gdef\csname CGTip@eq:o11\endcsname{Equation: delta\_X:=frac pq(e\textasciicircum{}c/X-1)}
\expandafter\gdef\csname CGTip@eq:perturbedfloor\endcsname{Equation: lfloor e\textasciicircum{}log(p/q)+c/X(qk+a)rfloor = pk+ lfloor(pa)/(q)+delta\_X(qk+a)rfloor.}
\expandafter\gdef\csname CGTip@eq:o12\endcsname{Equation: rle (pa)/(q)+delta\_X(qk+a)<r+1.}
\expandafter\gdef\csname CGTip@eq:blockinterval\endcsname{Equation: fracr-pa/qXp(e\textasciicircum{}c/X-1) -(a)/(qX) le (k)/(X) < fracr+1-pa/qXp(e\textasciicircum{}c/X-1) -(a)/(qX).}
\expandafter\gdef\csname CGTip@eq:o13\endcsname{Equation: qk+a, pk+lfloor(pa)/(q)rfloor+r.}
\expandafter\gdef\csname CGTip@eq:o14\endcsname{Equation: qr=pa-qlfloor(pa)/(q)rfloor =pabmod qin1,ldots,q-1,}
\expandafter\gdef\csname CGTip@eq:o15\endcsname{Equation: n=qk, p(e\textasciicircum{}c/X-1)k<1.}
\expandafter\gdef\csname CGTip@eq:o16\endcsname{Equation: R\_X(c)= min ( lfloorfrac Xqrfloor, lceilfrac1p(e\textasciicircum{}c/X-1)rceil-1 ),}
\expandafter\gdef\csname CGTip@eq:o17\endcsname{Equation: fracR\_X(c)Xlongrightarrow min (frac1q,frac1pc),}
\expandafter\gdef\csname CGTip@eq:o18\endcsname{Equation: (pX)/(q)(e\textasciicircum{}c/X-1) = 1+(q)/(2pX)+O(X\textasciicircum{}-2),}
\expandafter\gdef\csname CGTip@eq:o19\endcsname{Equation: frac1X C\_X (log m+frac cX) longrightarrow 0, c<0, lambda(m), 0le cle1/m, lambda(m)/(mc), c>1/m.}
\expandafter\gdef\csname CGTip@eq:o20\endcsname{Equation: C\_X (log m+frac cX)=lambda(m)X}
\expandafter\gdef\csname CGTip@eq:naturalmajority\endcsname{Equation: liminf\_Xtoinfty lambda(m)frac1X C\_X (log m+frac cX) ge 2min (1,frac1mc)-1>0.}
\expandafter\gdef\csname CGTip@eq:o21\endcsname{Equation: I\_0,X:=1le nle X:m(e\textasciicircum{}c/X-1)n<1.}
\expandafter\gdef\csname CGTip@eq:o22\endcsname{Equation: lambda(m)C\_X (log m+frac cX)ge2|I\_0,X|-X.}
\expandafter\gdef\csname CGTip@eq:o23\endcsname{Equation: frac|I\_0,X|Xtomin (1,frac1mc),}
\expandafter\gdef\csname CGTip@eq:o24\endcsname{Equation: (log c\_X)/(log X)longrightarrowtheta, 0letheta<1.}
\expandafter\gdef\csname CGTip@eq:logmajority\endcsname{Equation: liminf\_Xtoinfty lambda(pq) C\_X\textasciicircum{}log (logfrac pq+(c\_X)/(X)) ge frac2(1-theta)q-1.}
\expandafter\gdef\csname CGTip@eq:o25\endcsname{Equation: liminf\_Xtoinfty lambda(m)C\_X\textasciicircum{}log (log m+(c\_X)/(X)) ge1-2theta>0.}
\expandafter\gdef\csname CGTip@eq:o26\endcsname{Equation: p(e\textasciicircum{}c\_X/X-1)k<1.}
\expandafter\gdef\csname CGTip@eq:o27\endcsname{Equation: K\_X=X\textasciicircum{}1-theta+o(1).}
\expandafter\gdef\csname CGTip@eq:o28\endcsname{Equation: B\_X= frac1H\_Xsum\_kle K\_Xfrac1qk = (1-theta)/(q)+o(1).}
\expandafter\gdef\csname CGTip@eq:o29\endcsname{Equation: lambda(pq)C\_X\textasciicircum{}logge B\_X-(1-B\_X)=2B\_X-1.}
\expandafter\gdef\csname CGTip@eq:o32\endcsname{Equation: C\_X (log m+frac cX) = sum\_nle Xlambda(n)lambda(mn-1).}
\expandafter\gdef\csname CGTip@eq:oai\endcsname{Equation: |sum\_nle Xlambda(a\_1n+b\_1)lambda(a\_2n+b\_2)|le C fracX(log X)\textasciicircum{}kappa (Xge3, Xinmathbb R).}
\expandafter\gdef\csname CGTip@eq:macro_rate\endcsname{Equation: sum\_substackA\_X<kle B\_X alpha\_1k+beta\_1ge1, alpha\_2k+beta\_2ge1lambda(alpha\_1k+beta\_1)lambda(alpha\_2k+beta\_2)=O(fracX(log X)\textasciicircum{}kappa).}
\expandafter\gdef\csname CGTip@eq:profile_uncond\endcsname{Equation: frac1X C\_X (logfrac pq+frac cX)=Psi\_p,q(c)+O\_p,q,c((log X)\textasciicircum{}-kappa),}
\expandafter\gdef\csname CGTip@eq:left_uncond\endcsname{Equation: C\_X (log m+frac cX)=sum\_nle Xlambda(n)lambda(mn-1)=O(frac X(log X)\textasciicircum{}kappa).}
\expandafter\gdef\csname CGTip@eq:trivial_resonance\endcsname{Equation: C\_X (frac cX)=sum\_2le nle Xlambda(n-1)lambda(n),}
\expandafter\gdef\csname CGTip@eq:Ck\endcsname{Equation: C\textasciicircum{}(k)\_X(u):=sum\_substack1le nle X lfloor e\textasciicircum{}u\_inrfloorge1 forall ilambda(n)prod\_i=1\textasciicircum{}k-1lambda(lfloor e\textasciicircum{}u\_inrfloor), C\textasciicircum{}(k),log\_X(u):=frac1H\_Xsum\_substack1le nle X lfloor e\textasciicircum{}u\_inrfloorge1 forall ifraclambda(n)prod\_ilambda(lfloor e\textasciicircum{}u\_inrfloor)n.}
\expandafter\gdef\csname CGTip@eq:uk\endcsname{Equation: u\_i=log(p\_i)/(q\_i)+(c\_i)/(X), L:=lcm(q\_1,dots,q\_k-1), delta\_i:=(p\_i)/(q\_i)(e\textasciicircum{}c\_i/X-1).}
\expandafter\gdef\csname CGTip@eq:theta_k\endcsname{Equation: theta\_a,i:=(p\_ia)/(q\_i), b\_a,i:=lfloor(p\_ia)/(q\_i)rfloor, A\_i:=(p\_iL)/(q\_i)inmathbb N .}
\expandafter\gdef\csname CGTip@eq:floor_k\endcsname{Equation: lfloor e\textasciicircum{}u\_inrfloor=A\_im+b\_a,i+r\_i(n), r\_i(n):=lfloortheta\_a,i+delta\_inrfloor,}
\expandafter\gdef\csname CGTip@eq:forms_k\endcsname{Equation: f\_a,r,0(m):=Lm+a, f\_a,r,i(m):=A\_im+b\_a,i+r\_i (1le ile k-1).}
\expandafter\gdef\csname CGTip@eq:dec0_k\endcsname{Equation: C\textasciicircum{}(k)\_X(u)=sum\_a=0\textasciicircum{}L-1sum\_rin R sum\_min I\_a,r(X) prod\_i=0\textasciicircum{}k-1lambda(f\_a,r,i(m))}
\expandafter\gdef\csname CGTip@eq:dec_k\endcsname{Equation: C\textasciicircum{}(k)\_X(u)=sum\_substack(a,r) k'(a,r)=0varepsilon\_a,r |I\_a,r(X)| +sum\_substack(a,r) k'(a,r)ge1varepsilon\_a,rsum\_min I\_a,r(X) prod\_mathcal CinOdd(a,r)lambda(g\_mathcal C(m)).}
\expandafter\gdef\csname CGTip@eq:collapse\endcsname{Equation: prod\_i=0\textasciicircum{}k-1lambda(f\_a,r,i(m))=lambda(prod\_i d\_i)prod\_mathcal Clambda(g\_mathcal C(m))\textasciicircum{}|mathcal C|=varepsilon\_a,rprod\_mathcal CinOdd(a,r)lambda(g\_mathcal C(m)).}
\expandafter\gdef\csname CGTip@eq:profile_k\endcsname{Equation: Psi\textasciicircum{}(k)\_p,q(c):=frac1Lsum\_substack(a,r) k'(a,r)=0varepsilon\_a,r |J\_a,r(c)| .}
\expandafter\gdef\csname CGTip@eq:limit_k\endcsname{Equation: frac1XC\textasciicircum{}(k)\_X(u)longrightarrowPsi\textasciicircum{}(k)\_p,q(c) (Xtoinfty).}
\expandafter\gdef\csname CGTip@eq:ex4\endcsname{Equation: frac1X C\textasciicircum{}(4)\_X(log2+frac1X, log3+fracc\_2X, log6+frac1X)=frac13+O((log X)\textasciicircum{}-kappa).}
\expandafter\gdef\csname CGTip@eq:plateau_k\endcsname{Equation: C\textasciicircum{}(k)\_X(u)=lambda(m\_1cdots m\_k-1)sum\_nle Xlambda(n)\textasciicircum{}k= lambda(m\_1cdots m\_k-1)lfloor Xrfloor,  k even, [1mm] O(Xe\textasciicircum{}-c\_0sqrtlog X),  k odd.}
\expandafter\gdef\csname CGTip@eq:majority_k\endcsname{Equation: liminf\_Xtoinfty lambda(m\_1cdots m\_k-1)fracC\textasciicircum{}(k)\_X(u)X ge 2min(1,frac1max\_im\_ic\_i)-1 .}
\expandafter\gdef\csname CGTip@eq:left_k\endcsname{Equation: C\textasciicircum{}(k)\_X(u)=sum\_substacknle X m\_in-1ge1 forall ilambda(n)prod\_i=1\textasciicircum{}k-1lambda(m\_in-1).}
\expandafter\gdef\csname CGTip@eq:level3_23\endcsname{Equation: C\textasciicircum{}(3)\_X(log2+fracc\_1X, log3+fracc\_2X) =sum\_nle Xlambda(n)lambda(2n-1)lambda(3n-1)  =sum\_substack3le Mle 6X-3 Mequiv3 (mod 6)lambda(M)lambda(M+1)lambda(M+3),}
\expandafter\gdef\csname CGTip@eq:chowla013\endcsname{Equation: sum\_substackMle Y Mequiv3 (mod 6)lambda(M)lambda(M+1)lambda(M+3)=o(Y).}
\expandafter\gdef\csname CGTip@eq:Pr\endcsname{Equation: P\_r(Y):=sum\_mle Ylambda(m)prod\_i=1\textasciicircum{}k-1lambda(a\_im+r\_i), pi\_r:=varepsilon\_0,r mathbf 1[k'(0,r)=0],}
\expandafter\gdef\csname CGTip@eq:last_block\endcsname{Equation: I\_r(X)=[alpha\_rX, X]+O(1), alpha\_r:=max\_i(r\_i)/(r\_i+1)in[tfrac12,1).}
\expandafter\gdef\csname CGTip@eq:suffix\endcsname{Equation: P\_r(Y)-P\_r(alpha\_rY)=pi\_r(1-alpha\_r)Y+o(Y) (Ytoinfty).}
\expandafter\gdef\csname CGTip@eq:231\endcsname{Equation: sum\_mle Xlambda(m)lambda(2m+1)lambda(3m+1)=o(X):}
\expandafter\gdef\csname CGTip@eq:oai_params\endcsname{Equation: varepsilon=10\textasciicircum{}-4, eta=(varepsilon)/(100), kappa\_0=(400)/(eta), A=e\textasciicircum{}2kappa\_0, c\_*=(varepsilon)/(100), theta=frac14,}
\expandafter\gdef\csname CGTip@eq:oai_primes\endcsname{Equation:  mathcal C=p: B\textasciicircum{}1-eta<log ple B, mathcal Z=p: B\textasciicircum{}1-varepsilon<log ple B\textasciicircum{}1-eta,  mathcal S=mathcal Ccupmathcal Z, P\_0=e\textasciicircum{}B\textasciicircum{}1-varepsilon,}
\expandafter\gdef\csname CGTip@eq:oai124\endcsname{Equation: limsup\_Ytoinfty|frac1Ysum\_zle Yb\_1(z)sum\_M<wletau M(c\_w)/(w) b\_2(z+hw)|ll B\textasciicircum{}-1-varepsilon,}
\expandafter\gdef\csname CGTip@eq:oai127\endcsname{Equation: limsup\_Ytoinfty sup\_alphainmathbb R/mathbb Z frac1DYint\_0\textasciicircum{}Y|sum\_y<mle y+Db(m)e(alpha m)| dyll(loglog D)/(log D),}
\expandafter\gdef\csname CGTip@eq:E_k\endcsname{Equation: mathcal E\_k:=limsup\_Ytoinfty|frac1Ysum\_zle Yb\_0(z)sum\_M<wletau M(c\_w)/(w)prod\_i=1\textasciicircum{}k-1b\_i(z+h\_iw)|ll B\textasciicircum{}-1-varepsilon.}
\expandafter\gdef\csname CGTip@eq:cw_general\endcsname{Equation: lambda(u)\textasciicircum{}k c\textasciicircum{}(u)\_w=lambda(w)\textasciicircum{}k(-theta)\textasciicircum{}omega(w)mathbf 1[uwinmathcal D]= theta\textasciicircum{}omega(w)mathbf 1[uwinmathcal D], k odd, (-theta)\textasciicircum{}omega(w)mathbf 1[uwinmathcal D], k even.}
\expandafter\gdef\csname CGTip@eq:barrier_poly_127\endcsname{Equation: limsup\_Ytoinftyfrac1DYint\_0\textasciicircum{}Y|sum\_y<mle y+Db\_i(m)e(alpha m)| dyle D\textasciicircum{}-1/2;}
\expandafter\gdef\csname CGTip@eq:barrier_poly_E\endcsname{Equation: frac1Ysum\_zle Yb\_0(z)sum\_M<wletau M(c\_w)/(w)prod\_i=1\textasciicircum{}k-1b\_i(z+h\_iw)=fraclfloor YrfloorYsum\_substackM<wletau M w P\_0-roughfrac1w .}
\expandafter\gdef\csname CGTip@eq:barrier_cs\endcsname{Equation: frac1Yint\_0\textasciicircum{}Y|S\_y(alpha)|\textasciicircum{}2dy=sum\_1le t,t'le D e(alpha(t-t'))  timesfrac1Yint\_0\textasciicircum{}Ye(beta R\_i((n\_0+t)\textasciicircum{}k-1-(n\_0+t')\textasciicircum{}k-1))dy .}
\expandafter\gdef\csname CGTip@eq:q_choice\endcsname{Equation: qgemax(2h\_2+1, 7h\_2(D\textasciicircum{}*+1)\textasciicircum{}3, 6|rho|tau\textasciicircum{}2H\textasciicircum{}2).}
\expandafter\gdef\csname CGTip@eq:barrier_pos_127\endcsname{Equation: limsup\_Ytoinftyfrac1DYint\_0\textasciicircum{}Y|sum\_y<mle y+Db\_i(m)e(alpha m)|dyle(3+log q)/(D);}
\expandafter\gdef\csname CGTip@eq:barrier_pos\endcsname{Equation: liminf\_Ytoinfty Re frac1Ysum\_zle Yb\_0(z) sum\_M<wletau M(c\_w)/(w)b\_1(z+h\_1w)b\_2(z+h\_2w)  ge frac12sum\_M<wletau M(c\_w)/(w).}
\expandafter\gdef\csname CGTip@eq:gowers\endcsname{Equation: |g|\_U\textasciicircum{}s\textasciicircum{}2\textasciicircum{}s:=mathbb E\_a |Delta\_ag|\_U\textasciicircum{}s-1\textasciicircum{}2\textasciicircum{}s-1 (sge2),}
\expandafter\gdef\csname CGTip@eq:vonneumann\endcsname{Equation: |mathbb E\_z,winmathbb Z\_Nprod\_i=0\textasciicircum{}k-1g\_i(z+h\_iw)|le|g\_k-1|\_U\textasciicircum{}k-1.}
\expandafter\gdef\csname CGTip@eq:vn_step\endcsname{Equation: |mathbb E\_z,wprod\_i=0\textasciicircum{}k-1g\_i(z+h\_iw)|\textasciicircum{}2 lemathbb E\_z|mathbb E\_wprod\_i=1\textasciicircum{}k-1g\_i(z+h\_iw)|\textasciicircum{}2  =mathbb E\_t mathbb E\_z,wprod\_i=1\textasciicircum{}k-1(Delta\_h\_itg\_i)(z+(h\_i-h\_1)w),}
\expandafter\gdef\csname CGTip@eq:local_control\endcsname{Equation:  frac1Y|sum\_zle Yb\_0(z)sum\_winmathcal Jfracc\_wwprod\_i=1\textasciicircum{}k-1b\_i(z+h\_iw)|   le C (varpilog(3M))\textasciicircum{}1/2 frac1Yint\_0\textasciicircum{}Y|b\_k-1|\_U\textasciicircum{}k(y;D) dy   +((2D+1)varpi)/(Y).}
\expandafter\gdef\csname CGTip@eq:windows\endcsname{Equation: |sum\_zle Yg(z)|lefrac1Dint\_0\textasciicircum{}Y|S\_y| dy+(2D+1)varpi .}
\expandafter\gdef\csname CGTip@eq:cs_w\endcsname{Equation: |S\_y|\textasciicircum{}2levarpisum\_winmathcal J(|c\_w|)/(w)|X\_y,w|\textasciicircum{}2le(varpi)/(M)sum\_winmathcal J|X\_y,w|\textasciicircum{}2 .}
\expandafter\gdef\csname CGTip@eq:Tk\endcsname{Equation: sum\_winmathcal J|X\_y,w|\textasciicircum{}2=sum\_|s|<DT\_s, T\_s:=sum\_xsum\_winmathcal J(Delta\_sF\_0)(x)prod\_i=1\textasciicircum{}k-1(Delta\_sG\_i)(x+h\_iw),}
\expandafter\gdef\csname CGTip@eq:Tk_bound\endcsname{Equation: |T\_s|le N\textasciicircum{}2(2+log N) |Delta\_sG\_k-1|\_U\textasciicircum{}k-1(mathbb Z\_N) .}
\expandafter\gdef\csname CGTip@eq:sum_k\endcsname{Equation: sum\_|s|<D|Delta\_sG\_k-1|\_U\textasciicircum{}k-1le N\textasciicircum{}1-2\textasciicircum{}1-k(sum\_sinmathbb Z\_N|Delta\_sG\_k-1|\textasciicircum{}2\textasciicircum{}k-1\_U\textasciicircum{}k-1)\textasciicircum{}2\textasciicircum{}1-k=N |G\_k-1|\textasciicircum{}2\_U\textasciicircum{}k.}
\expandafter\gdef\csname CGTip@eq:LU\endcsname{Equation: limsup\_Ytoinftyfrac1Yint\_0\textasciicircum{}Y|b|\_U\textasciicircum{}k(y;D) dyle K(2+sum\_pinmathcal Pfrac1p)\textasciicircum{}K(log D)\textasciicircum{}-A.}
\expandafter\gdef\csname CGTip@eq:Ek_bound\endcsname{Equation: mathcal E\_kll(B\textasciicircum{}-1+varepsilon(log B) B)\textasciicircum{}1/2(2+log B)\textasciicircum{}K B\textasciicircum{}-A(1-varepsilon)ll B\textasciicircum{}varepsilon/2(log B)\textasciicircum{}K+1/2B\textasciicircum{}-A(1-varepsilon),}
\expandafter\gdef\csname CGTip@eq:gk\endcsname{Equation: g\textasciicircum{}(k)\_d(x):=A\textasciicircum{}omega\_mathcal C(d) mathbf 1\_d\_mathcal Cmid xprod\_substackpmid d pinmathcal Z(mathbf 1\_pmid x-fractheta p)prod\_i=0\textasciicircum{}k-1w\_d(x+h\_id).}
\expandafter\gdef\csname CGTip@eq:Gk\endcsname{Equation: limsup\_Xtoinftyfrac1X|sum\_xle Xsum\_dinmathcal Da\_dprod\_i=0\textasciicircum{}k-1F\_i(x+h\_id) g\textasciicircum{}(k)\_d(x)|ll L\_0B\textasciicircum{}-1-c\_*/2,}
\expandafter\gdef\csname CGTip@eq:conditional_k\endcsname{Equation: frac1Nsum\_substackmle N mequiv b (mod l)prod\_i=0\textasciicircum{}k-1lambda(m+h\_i)longrightarrow0 (Ntoinfty).}
\expandafter\gdef\csname CGTip@eq:char_div\endcsname{Equation: |mathbb E\_divchi(d)| =prod\_pinmathcal S|(1+A\_achi(p)/p)/(1+A\_a/p)|  leexp(-sum\_pinmathcal Z(1-Re chi(p))/(p)+O(sum\_p>P\_0p\textasciicircum{}-2))to0,}
\expandafter\gdef\csname CGTip@eq:RB_k\endcsname{Equation: R\_B(X):=sum\_substackxle X xequiv b (l)sum\_dinmathcal Da\_dA\textasciicircum{}omega\_mathcal C(d)prod\_i=0\textasciicircum{}k-1lambda(x+h\_id)w\_d(x+h\_id) mathbf 1\_dmid x.}
\expandafter\gdef\csname CGTip@eq:channels\endcsname{Equation: C\_H(X)=mu\_p(H)T\_p,H,0(X)+sum\_j=1\textasciicircum{}Jsum\_substackubmod p\textasciicircum{}j pnmid uwidehatDelta\_p,H,j(u) T\_p,H,j,u(X)+O(kp\textasciicircum{}-J+frac kX),}
\expandafter\gdef\csname CGTip@eq:o37\endcsname{Equation: c   lambda(2)C\_X/X   (2|I\_0,X|-X)/X   2min(1,(2c)\textasciicircum{}-1)-1 hline 0.60   0.833072800   0.666666600   0.666666667 0.90   0.555258200   0.111111000   0.111111111}
\expandafter\gdef\csname CGTip@sec:introduction\endcsname{Section: Introduction}
\expandafter\gdef\csname CGTip@sec:input\endcsname{Section: Affine input and logarithmic normalization}
\expandafter\gdef\csname CGTip@sec:rational\endcsname{Section: Exact rational decomposition}
\expandafter\gdef\csname CGTip@sec:logstep\endcsname{Section: Unconditional logarithmic microscopic step}
\expandafter\gdef\csname CGTip@sec:natural\endcsname{Section: Conditional natural microscopic profile}
\expandafter\gdef\csname CGTip@sec:integer\endcsname{Section: Integer rigidity and unconditional majority}
\expandafter\gdef\csname CGTip@sec:unconditional\endcsname{Section: Unconditional natural profile}
\expandafter\gdef\csname CGTip@sec:kpoint\endcsname{Section: The k-point resonance family}
\expandafter\gdef\csname CGTip@subsec:kpoint_def\endcsname{Subsection: Definition and exact decomposition}
\expandafter\gdef\csname CGTip@subsec:kpoint_suff\endcsname{Subsection: What each order of Chowla closes}
\expandafter\gdef\csname CGTip@subsec:kpoint_integer\endcsname{Subsection: Plateaus, majority and the left-hand identity at integer resonances}
\expandafter\gdef\csname CGTip@subsec:kpoint_nec\endcsname{Subsection: Necessity}
\expandafter\gdef\csname CGTip@sec:three\endcsname{Section: Three and more points: the centering step of the bias method}
\expandafter\gdef\csname CGTip@subsec:bias\endcsname{Subsection: The bias method for two points}
\expandafter\gdef\csname CGTip@subsec:barrier\endcsname{Subsection: Barriers}
\expandafter\gdef\csname CGTip@subsec:local_gowers\endcsname{Subsection: Local Gowers norms control the centering error}
\expandafter\gdef\csname CGTip@subsec:conditional\endcsname{Subsection: A conditional extension of the bias method}
\expandafter\gdef\csname CGTip@sec:numerics\endcsname{Section: Numerical verification}
\expandafter\gdef\csname CGTip@sec:conclusions\endcsname{Section: Conclusions}
\expandafter\gdef\csname CGTip@sec:data\endcsname{Section: Data and reproducibility}
\expandafter\gdef\csname CGTip@sec:ai\endcsname{Section: Use of AI tools}
\expandafter\gdef\csname CGTip@fig:rational\endcsname{Figure: Measured rational resonances at X=2cdot10\textasciicircum{}6 and the conditional natural main term.}
\expandafter\gdef\csname CGTip@fig:profile\endcsname{Figure: Natural microscopic profile around log(3/2).}
\expandafter\gdef\csname CGTip@fig:logmajority\endcsname{Figure: Finite logarithmic mesoscopic check for m=2, X=10\textasciicircum{}7. The dashed curve is the envelope supplied by the unconditional signed majority bound.}
\makeatother

\hypersetup{
  colorlinks=false,
  pdfborder={0 0 1},
  linkbordercolor={1 0 0},
  citebordercolor={0 0.3 1},
  urlbordercolor={0 0.6 0.6},
  filebordercolor={0 0.6 0.6},
  linktoc=all,
  bookmarksnumbered=true,
  bookmarksopen=true,
  bookmarksopenlevel=1,
  pdfstartview=FitH,
  pdftitle={Rational Resonances and Microscopic Natural--Logarithmic Profiles for the Liouville Function (expanded version)},
  pdfauthor={V. Casa\~na-Giner},
  pdfsubject={Liouville function, Chowla conjecture, rational resonances, k-point family},
  pdfkeywords={Liouville function; Chowla conjecture; rational resonance; microscopic profile; local Gowers uniformity; divisor graph}
}
% Back-references in the bibliography.
\renewcommand*{\backref}[1]{}
\renewcommand*{\backrefalt}[4]{%
  \ifcase #1 %
    \relax
  \or
    \ (Cited on p.~#2.)%
  \else
    \ (Cited on pp.~#2.)%
  \fi}

\numberwithin{equation}{section}
\setcounter{secnumdepth}{3}
\setcounter{tocdepth}{3}
\emergencystretch=1.5em

\theoremstyle{plain}
\newtheorem{theorem}{Theorem}[section]
\newtheorem{proposition}[theorem]{Proposition}
\newtheorem{lemma}[theorem]{Lemma}
\newtheorem{corollary}[theorem]{Corollary}
\newtheorem{conjecture}[theorem]{Conjecture}

\theoremstyle{definition}
\newtheorem{definition}[theorem]{Definition}
\newtheorem{hypothesis}[theorem]{Hypothesis}
\newtheorem{example}[theorem]{Example}

\theoremstyle{definition}
\newtheorem{experiment}[theorem]{Experiment}

\theoremstyle{remark}
\newtheorem{remark}[theorem]{Remark}

\begin{document}

\begin{frontmatter}

\title{Rational Resonances and Microscopic Natural--Logarithmic Profiles for the Liouville Function}

\author[uned]{Víctor Casaña-Giner\texorpdfstring{\corref{cor1}}{}}
\ead{vcasana1@alumno.uned.es}
\ead[url]{https://orcid.org/0000-0002-8252-2608}
\cortext[cor1]{Corresponding author.}
\affiliation[uned]{organization={Universidad Nacional de Educación a Distancia},
addressline={C/ Bravo Murillo, 38}, postcode={28015}, city={Madrid}, country={Spain}}

\begin{abstract}
Let $\lambda(n)=(-1)^{\Omega(n)}$. For $X\in\mathbb N$ and $u\in\mathbb R$, define
\begin{equation}
C_X(u):=\sum_{\substack{1\le n\le X\\ \lfloor e^u n\rfloor\ge1}}
\lambda(n)\lambda(\lfloor e^u n\rfloor).
\label{eq:o1}
\end{equation}
For reduced rational dilations $e^u=p/q$, we isolate exactly the unique
proportional residue class. Assuming the standard fixed two-point affine
Chowla conjecture, this yields the natural microscopic profile
\begin{equation}
\frac{1}{X}C_X\!\left(\log\frac pq+\frac cX\right)
\longrightarrow
\begin{cases}
0,&c<0,\\
\lambda(pq)/q,&0\le c\le q/p,\\
\lambda(pq)/(pc),&c>q/p.
\end{cases}
\label{eq:o2}
\end{equation}
The logarithmically weighted analogue has an unconditional one-sided
microscopic limit at the same $X^{-1}$ scale. At integer resonances we prove,
without Chowla input, a right-hand natural-average majority window
$0<c<2/m$ and a logarithmic mesoscopic majority theorem whose elementary
majority mechanism has threshold exponent $1/2$: every sequence
$c_X=X^{\theta+o(1)}$ with $\theta<1/2$ forces non-cancellation at the
corresponding right-hand scale $c_X/X=X^{-1+\theta+o(1)}$. Thus natural and logarithmic weighting exhibit different
microscopic scales, while integer resonances are singled out by unconditional
majority rigidity.
In this expanded version the two-point input is supplied by the ordinary
two-point theorem of OpenAI (2026), which makes the natural profile
unconditional with the rate $(\log X)^{-\kappa}$. We introduce the $k$-point
resonance family, prove its exact decomposition into proportionality classes
and its closed-form profile, extend the plateaus, the majority theorem and the
left-hand identity to every $k$, and show that the level-three family at the
resonances $(2,3)$ is, for every large $X$, the ordinary three-point sum
$\sum_{M\equiv3\,(6)}\lambda(M)\lambda(M+1)\lambda(M+3)$. For the bias method
behind the two-point theorem we prove that its centering step cannot be
extended to $k\ge3$ with the short-interval input it uses, that local $U^k$
uniformity suffices for that step, and that a quantitative local Gowers
uniformity hypothesis together with a $k$-linear divisor-graph hypothesis
yields $k$-point Chowla in progressions and the profile of the whole family.
\end{abstract}

\begin{keyword}
Liouville function \sep Chowla conjecture \sep Beatty sequence \sep
rational resonance \sep logarithmic average \sep microscopic profile
\end{keyword}

\end{frontmatter}

\tableofcontents

\section{Introduction}
\label{sec:introduction}

The Liouville function $\lambda(n)=(-1)^{\Omega(n)}$ is completely
multiplicative. The fixed two-point Chowla conjecture predicts natural-average
cancellation of $\lambda$ on non-proportional affine forms. Tao proved the
corresponding logarithmically averaged two-point theorem
\cite{Tao2016}, while Matom\"aki--Radziwi{\l}{\l}--Tao proved an averaged
natural form \cite{MRT2015}. Arithmetic functions on Beatty sequences were
studied by Abercrombie \cite{Abercrombie1995} and
Abercrombie--Banks--Shparlinski \cite{AbercrombieBanksShparlinski2009}.

For irrational $\alpha$ with first slope equal to $1$, logarithmic
cancellation of
$\lambda(n)\lambda(\lfloor\alpha n\rfloor)$ was proved by
Crn\v{c}evi\'c, Hern\'andez, Rizk, Sereesuchart and R.~Tao
\cite{CHRST2024}. For a general irrational ratio of slopes, logarithmic
two-point cancellation is due to Ter\"av\"ainen and Walker
\cite[Theorem~1.2]{TeravainenWalker2025}; they moreover prove that in the
rationally dependent case the logarithmic limit exists and characterize when
it vanishes, answering the two-point case of a conjecture of Frantzikinakis
\cite{Frantzikinakis2020}. The statement of Ter\"av\"ainen--Walker, Theorem~1.2(3), gives existence
and a zero/non-zero criterion in the rationally dependent case, but it does
not state a closed-form value for the present specialization. Here
$\alpha_1=1$, $\alpha_2=p/q$ and $r=q/p$. Their vanishing criterion cannot
hold eventually, because for every $\ell$ divisible by $q$,
\begin{equation}
\ell=\frac qp\left\lfloor\frac{p\ell}{q}\right\rfloor.
\label{eq:o3}
\end{equation}
Thus their theorem implies non-vanishing in this case. The explicit value
$\lambda(pq)/q$ is obtained here by isolating the proportional residue class
and applying Tao's logarithmic two-point theorem to the finitely many
non-proportional affine classes. The new object is the singular
$X$-dependent perturbation
\begin{equation}
u_X=\log(p/q)+c/X,
\label{eq:o4}
\end{equation}
together with its natural microscopic profile and the unconditional majority
phenomena that appear at integer resonances.

The structure is transparent. At a reduced rational dilation $p/q$, the
residue class $n\equiv0\pmod q$ is the unique proportional affine class; all
other residue classes are non-proportional. The natural main term is therefore
the sign $\lambda(pq)$ times the size of this proportional block, while the
remaining pieces require affine Chowla cancellation.

\paragraph{Expanded version}
Sections~\CGref{sec:unconditional}--\CGref{sec:three} are new. The remainder of
the residue-class decomposition is a finite set of fixed binary affine
correlations, and the ordinary two-point Chowla theorem of
\cite[Theorem~1.1, p.~3]{OpenAI2026} controls exactly these; Section~\CGref{sec:unconditional}
records the resulting unconditional profile with a rate, and shows that at the
trivial resonance $p=q=1$, $-1<c<0$, the profile is literally the shift-one
Chowla sum. Section~\CGref{sec:kpoint} constructs the $k$-point family
$C^{(k)}_X$, proves its exact decomposition (Theorem~\CGref{thm:dec_k}), the
criterion for unconditional closure (Proposition~\CGref{prop:criterion}), the
plateau, majority and left-hand identities at every level
(Theorem~\CGref{thm:integer_k}), the identification of level three at $(2,3)$
with an ordinary three-point Chowla sum (Corollary~\CGref{cor:level3_23}), and
the necessity theorem (Theorem~\CGref{thm:nec_k}). Section~\CGref{sec:three}
examines the bias method of \cite[Part~II]{OpenAI2026} for $k\ge3$ points:
a barrier for its centering step (Propositions~\CGref{prop:barrier_poly}
and~\CGref{prop:barrier_pos}), an unconditional bound of the $k$-point centering
error by local Gowers norms (Proposition~\CGref{prop:local_control}), and a
conditional $k$-point theorem (Theorem~\CGref{thm:conditional_k},
Corollary~\CGref{cor:family_conditional}) whose two hypotheses are stated
explicitly and are not proved here.

\section{Affine input and logarithmic normalization}
\label{sec:input}

\begin{conjecture}[Standard fixed two-point affine Chowla]
\label{conj:affine_chowla}
Let $a_1,a_2\in\mathbb N$ and $b_1,b_2\in\mathbb Z$ be fixed, with the two
affine forms positive for all sufficiently large $n$ and
$a_1b_2-a_2b_1\ne0$. Then
\begin{equation}
\sum_{\substack{n\le Y\\ a_1n+b_1\ge1,\ a_2n+b_2\ge1}}\lambda(a_1n+b_1)\lambda(a_2n+b_2)=o(Y).
\label{eq:o5}
\end{equation}
\end{conjecture}

No uniform or quantitative strengthening of this conjecture is used.

Throughout the paper, $X\in\mathbb N$ and all limits $X\to\infty$ are taken
through positive integers. For $u\in\mathbb R$ we write
\begin{equation}
C_X(u):=
\sum_{\substack{1\le n\le X\\ \lfloor e^u n\rfloor\ge1}}
\lambda(n)\lambda(\lfloor e^u n\rfloor).
\label{eq:CXdef}
\end{equation}

\begin{lemma}[Macroscopic subintervals]
\label{lem:macroscopic}
Assume Conjecture~\CGref{conj:affine_chowla} for a fixed pair of affine forms.
If $0\le A_X\le B_X$, $A_X/X\to A$ and $B_X/X\to B$ with $0\le A\le B<\infty$, then
\[
\sum_{\substack{A_X<n\le B_X\\ a_1n+b_1\ge1,\ a_2n+b_2\ge1}}\lambda(a_1n+b_1)\lambda(a_2n+b_2)=o(X).
\]
\end{lemma}

\begin{proof}
Let $S(Y)$ be the sum in \CGeqref{eq:o5}. The conjecture says that for every
$\varepsilon>0$ there is $C_\varepsilon$ with $|S(Y)|\le\varepsilon Y+C_\varepsilon$ for all
$Y\ge0$. The sum in the lemma is $S(B_X)-S(A_X)$, so it is at most
$\varepsilon(A_X+B_X)+2C_\varepsilon\le3\varepsilon BX$ for large $X$ (and
$\le\varepsilon X+2C_\varepsilon$ if $B=0$). This covers $A=0$ and $A=B$.
\end{proof}

Let
\begin{equation}
H_X:=\sum_{n\le X}\frac1n
\label{eq:o6}
\end{equation}
and define
\begin{equation}
C_X^{\log}(u):=
\frac1{H_X}
\sum_{\substack{1\le n\le X\\ \lfloor e^u n\rfloor\ge1}}
\frac{\lambda(n)\lambda(\lfloor e^u n\rfloor)}n.
\label{eq:logdef}
\end{equation}
Thus $|C_X^{\log}(u)|\le1$.

\begin{lemma}[Affine logarithmic cancellation with residue-class weight]
\label{lem:weighted_tao}
Let $A_1,A_2,Q\in\mathbb N$ and $B_1,B_2,A\in\mathbb Z$ be fixed.
Assume that the affine forms $A_1k+B_1$ and $A_2k+B_2$ are positive for all
sufficiently large $k$ and are non-proportional:
\begin{equation}
A_1B_2-A_2B_1\ne0.
\label{eq:weighted_tao_nonprop}
\end{equation}
If $Qk+A>0$ for all sufficiently large $k$, then
\begin{equation}
\sum_{\substack{1\le k\le K\\A_1k+B_1\ge1,\ A_2k+B_2\ge1\\Qk+A\ge1}}
\frac{\lambda(A_1k+B_1)\lambda(A_2k+B_2)}{Qk+A}
=o(\log K).
\label{eq:weighted_tao}
\end{equation}
\end{lemma}

\begin{proof}
Tao's logarithmically averaged two-point theorem
\cite[Theorem~1.2]{Tao2016} states that for fixed
$a_1,a_2\in\mathbb N$ and $b_1,b_2\in\mathbb Z$ with
$a_1b_2-a_2b_1\ne0$, and for any
$1\le\omega(K)\le K$ with $\omega(K)\to\infty$,
\begin{equation}
\sum_{K/\omega(K)<k\le K}
\frac{\lambda(a_1k+b_1)\lambda(a_2k+b_2)}{k}
=o(\log\omega(K)).
\label{eq:tao_theorem_form}
\end{equation}
Apply this with $(a_i,b_i)=(A_i,B_i)$ and $\omega(K)=K$.
Deleting or restoring finitely many initial terms changes the sum by $O(1)$,
so
\begin{equation}
\sum_{k\le K}
\frac{\lambda(A_1k+B_1)\lambda(A_2k+B_2)}{k}
=o(\log K).
\label{eq:tao_weight_one_over_k}
\end{equation}
The residue-class weight satisfies
\begin{equation}
\frac1{Qk+A}=\frac1Q\,\frac1k+O_{Q,A}(k^{-2}).
\label{eq:residue_weight_expansion}
\end{equation}
Since the Liouville product has modulus one, the total contribution of the
$O(k^{-2})$ error is $O(1)$. Combining
\CGeqref{eq:tao_weight_one_over_k} and
\CGeqref{eq:residue_weight_expansion} proves \CGeqref{eq:weighted_tao}.
\end{proof}

\begin{remark}[Status of Conjecture~\CGref{conj:affine_chowla}]
\label{rem:status_conj}
For non-negative intercepts, Conjecture~\CGref{conj:affine_chowla} is now
\cite[Theorem~1.1, p.~3]{OpenAI2026}, an unrefereed preprint, with a power-of-logarithm saving at
every real cutoff; negative intercepts reduce to that case by a fixed shift
(Lemma~\CGref{lem:macro_rate}). Sections~\CGref{sec:rational}--\CGref{sec:integer}
are kept in their original conditional formulation; the unconditional
statements are collected in Section~\CGref{sec:unconditional}.
\end{remark}

\section{Exact rational decomposition}
\label{sec:rational}

\begin{proposition}[Exact residue-class decomposition]
\label{prop:decomposition}
Let $(p,q)=1$ and $X\in\mathbb N$. Then
\begin{align}
C_X\!\left(\log\frac pq\right)
&=
\lambda(pq)\left\lfloor\frac Xq\right\rfloor \nonumber\\
&\quad+
\sum_{a=1}^{q-1}
\sum_{\substack{k\ge0\\1\le qk+a\le X\\
pk+\lfloor pa/q\rfloor\ge1}}
\lambda(qk+a)
\lambda\!\left(pk+\left\lfloor\frac{pa}{q}\right\rfloor\right).
\label{eq:decomp}
\end{align}
For every $a\ne0$ the two affine forms in the inner sum are
non-proportional.
\end{proposition}

\begin{proof}
Write $n=qk+a$, $0\le a<q$. Then
\begin{equation}
\left\lfloor\frac pq(qk+a)\right\rfloor
=
pk+\left\lfloor\frac{pa}{q}\right\rfloor.
\label{eq:o7}
\end{equation}
For $a=0$ and $k\ge1$, complete multiplicativity gives
$\lambda(qk)\lambda(pk)=\lambda(pq)$. For $a\ne0$, proportionality would
require $q\lfloor pa/q\rfloor=pa$, impossible because $(p,q)=1$ and
$q\nmid a$.
\end{proof}

\begin{theorem}[Fixed rational natural resonance]
\label{thm:fixed}
Assume Conjecture~\CGref{conj:affine_chowla}. Then
\begin{equation}
\frac1X C_X\!\left(\log\frac pq\right)\longrightarrow\frac{\lambda(pq)}q.
\label{eq:o8}
\end{equation}
For $q=1$ the exact identity
\begin{equation}
C_X(\log p)=\lambda(p)X
\label{eq:o9}
\end{equation}
holds for every $X\in\mathbb N$.
\end{theorem}

\begin{proof}
The proportional class in Proposition~\CGref{prop:decomposition} contributes
$\lambda(pq)\lfloor X/q\rfloor$. Every other residue class is a fixed
non-proportional affine correlation and is $o(X)$ by
Conjecture~\CGref{conj:affine_chowla}. For $q=1$ there are no remaining
classes.
\end{proof}

\section{Unconditional logarithmic microscopic step}
\label{sec:logstep}

\begin{theorem}[Logarithmic microscopic step]
\label{thm:logstep}
Let $(p,q)=1$ and let $c\in\mathbb R$ be fixed. Then
\begin{equation}
C_X^{\log}\!\left(\log\frac pq+\frac cX\right)
\longrightarrow
\frac{\lambda(pq)}q\,\mathbf 1_{\{c\ge0\}}
\label{eq:logstep}
\end{equation}
unconditionally.
\end{theorem}

\begin{proof}
Fix $\varepsilon\in(0,1)$ and set
\begin{equation}
Y=X^{1-\varepsilon}.
\label{eq:logstep_Y}
\end{equation}
Write $n=qk+a$, $0\le a<q$. Uniformly for $n\le Y$,
\begin{equation}
\frac pq(e^{c/X}-1)n=O_{p,q,c}(X^{-\varepsilon}).
\label{eq:logstep_small_perturbation}
\end{equation}
If $a=0$, then $\{pa/q\}=0$. When $q=1$ there are no nonzero residue classes. If $a\ne0$, coprimality gives
\begin{equation}
\left\{\frac{pa}{q}\right\}\in
\left[\frac1q,1-\frac1q\right].
\label{eq:logstep_fractional_separation}
\end{equation}
Hence, for all sufficiently large $X$, the perturbation in
\CGeqref{eq:logstep_small_perturbation} cannot change the floor on any
$a\ne0$ residue class.

Suppose first that $c\ge0$. On $n\le Y$ the perturbed correlation agrees
exactly with the fixed rational dilation. The proportional class $a=0$
contributes
\begin{equation}
\frac{\lambda(pq)}{H_X}
\sum_{k\le Y/q}\frac1{qk}
=
\frac{H_Y}{H_X}\frac{\lambda(pq)}q+o(1).
\label{eq:logstep_prop_initial}
\end{equation}
For each $a\ne0$, the two fixed affine forms are non-proportional, so
Lemma~\CGref{lem:weighted_tao}, with weight $(qk+a)^{-1}$, shows that its
entire contribution on $n\le Y$ is $o(H_Y)$ and therefore $o(1)$ after
division by $H_X$. The tail $Y<n\le X$ has normalized harmonic mass
\begin{equation}
\frac{H_X-H_Y}{H_X}=\varepsilon+o(1).
\label{eq:logstep_tail_mass}
\end{equation}
Consequently,
\begin{equation}
\left|
C_X^{\log}\!\left(\log\frac pq+\frac cX\right)
-\frac{H_Y}{H_X}\frac{\lambda(pq)}q
\right|
\le
\frac{H_X-H_Y}{H_X}+o(1).
\label{eq:logstep_bilateral}
\end{equation}
Since $H_Y/H_X\to1-\varepsilon$, first letting $X\to\infty$ and then
$\varepsilon\downarrow0$ gives the branch $c\ge0$.
Ter\"av\"ainen--Walker's Theorem~1.2(3) already gives the rationally
dependent zero/non-zero alternative and implies non-vanishing in the present
homogeneous rational case; the closed-form value $\lambda(pq)/q$ is computed
here from the proportional residue class
\cite[Theorem~1.2(3)]{TeravainenWalker2025}.

Now suppose that $c<0$. On the class $a=0$ and the range $n=qk\le Y$,
\CGeqref{eq:logstep_small_perturbation} gives
\begin{equation}
\left\lfloor pe^{c/X}k\right\rfloor=pk-1
\label{eq:logstep_negative_zero_class}
\end{equation}
for every $k\ge1$ once $X$ is large, except for the possible initial term
for which the target equals zero and is excluded by definition. To apply
Tao in the usual non-negative-intercept formulation, write $k=\ell+1$:
\begin{equation}
qk=q\ell+q,\qquad pk-1=p\ell+(p-1).
\label{eq:logstep_reindex}
\end{equation}
If $p=1$, the single term with second value zero is omitted; the remaining
terms may be reindexed once more. The two forms are non-proportional. Every
$a\ne0$ class remains non-proportional by
\CGeqref{eq:logstep_fractional_separation}. Lemma~\CGref{lem:weighted_tao}
therefore makes the whole initial range $o(H_Y)$, while the tail is bounded
by \CGeqref{eq:logstep_tail_mass}. Letting $X\to\infty$ and then
$\varepsilon\downarrow0$ gives zero.
\end{proof}

\section{Conditional natural microscopic profile}
\label{sec:natural}

\begin{theorem}[Rational natural microscopic profile]
\label{thm:profile}
Assume Conjecture~\CGref{conj:affine_chowla}.\par
Let $(p,q)=1$ and fix $c\in\mathbb R$. Then
\begin{equation}
\frac1X C_X\!\left(\log\frac pq+\frac cX\right)
\longrightarrow
\Psi_{p,q}(c),
\label{eq:profilelimit}
\end{equation}
where
\begin{equation}
\Psi_{p,q}(c)=
\begin{cases}
0,&c<0,\\[1mm]
\dfrac{\lambda(pq)}q,&0\le c\le q/p,\\[3mm]
\dfrac{\lambda(pq)}{pc},&c>q/p.
\end{cases}
\label{eq:profile}
\end{equation}
For $c>0$,
\begin{equation}
\Psi_{p,q}(c)=
\lambda(pq)\min\!\left(\frac1q,\frac1{pc}\right).
\label{eq:o10}
\end{equation}
\end{theorem}

\begin{proof}
Put
\begin{equation}
\delta_X:=\frac pq(e^{c/X}-1)
\label{eq:o11}
\end{equation}
and write $n=qk+a$, $0\le a<q$. Then
\begin{equation}
\left\lfloor e^{\log(p/q)+c/X}(qk+a)\right\rfloor
=
pk+
\left\lfloor\frac{pa}{q}+\delta_X(qk+a)\right\rfloor.
\label{eq:perturbedfloor}
\end{equation}

Assume first $c>0$. For an integer $r$, the block on which the final floor in
\CGeqref{eq:perturbedfloor} equals $\lfloor pa/q\rfloor+r$ is
\begin{equation}
r\le
\left\{\frac{pa}{q}\right\}+\delta_X(qk+a)<r+1.
\label{eq:o12}
\end{equation}
Equivalently,
\begin{align}
\frac{r-\{pa/q\}}{Xp(e^{c/X}-1)}
-\frac{a}{qX}
\le \frac{k}{X}
<
\frac{r+1-\{pa/q\}}{Xp(e^{c/X}-1)}
-\frac{a}{qX}.
\label{eq:blockinterval}
\end{align}
This is written for $c>0$. For $c<0$ the factor $Xp(e^{c/X}-1)$ is negative, dividing by it reverses the inequalities, and \CGeqref{eq:blockinterval} holds with its two endpoints interchanged; nothing below depends on their order.
Since $Xp(e^{c/X}-1)\to pc$, both endpoints converge. After intersection
with $0\le k/X\le1/q+O(X^{-1})$, each block has endpoints converging to
fixed limits. If the limiting endpoints are distinct, the block is
macroscopic; if they coincide, its length is $o(X)$; otherwise it is empty.
For fixed $p,q,c$, the perturbation in \CGeqref{eq:perturbedfloor} is uniformly
bounded for $0\le n\le X$, so the number of integers $r$ that occur is
bounded in terms of $p,q,c$ and is independent of $X$.

On such a block the two affine forms are
\begin{equation}
qk+a,\qquad
pk+\left\lfloor\frac{pa}{q}\right\rfloor+r.
\label{eq:o13}
\end{equation}
For $a\ne0$, proportionality would require
\begin{equation}
qr=pa-q\left\lfloor\frac{pa}{q}\right\rfloor
=pa\bmod q\in\{1,\ldots,q-1\},
\label{eq:o14}
\end{equation}
which is impossible. For $a=0$, proportionality occurs only for $r=0$.
Thus every block except $(a,r)=(0,0)$ is $o(X)$ by
Lemma~\CGref{lem:macroscopic}.

The unique proportional block is
\begin{equation}
n=qk,\qquad p(e^{c/X}-1)k<1.
\label{eq:o15}
\end{equation}
Its cardinality is
\begin{equation}
R_X(c)=
\min\!\left(
\left\lfloor\frac Xq\right\rfloor,
\left\lceil\frac1{p(e^{c/X}-1)}\right\rceil-1
\right),
\label{eq:o16}
\end{equation}
up to the harmless endpoint convention. Hence
\begin{equation}
\frac{R_X(c)}X\longrightarrow
\min\!\left(\frac1q,\frac1{pc}\right),
\label{eq:o17}
\end{equation}
and every term in this block equals $\lambda(pq)$.

If $c<0$, the same block decomposition has negative corrections. For
$a=0$ the correction is a nonzero negative integer for every $k\ge1$, so no
proportional block remains. For $a\ne0$, proportionality would again imply
$qr=pa\bmod q$, impossible. Only finitely many negative correction integers occur for fixed $p,q,c$.
For each such correction, the target affine form can fail to be positive only
for $O_{p,q,c}(1)$ initial values of $k$; these terms are omitted by the
definition of $C_X$. All remaining blocks are $o(X)$.

At $c=0$, Theorem~\CGref{thm:fixed} applies. At the boundary $pc=q$,
\begin{equation}
\frac{pX}{q}(e^{c/X}-1)
=
1+\frac{q}{2pX}+O(X^{-2}),
\label{eq:o18}
\end{equation}
so only $O(1)$ terminal terms leave the proportional block. The normalized
limit remains $\lambda(pq)/q$.
\end{proof}

\section{Integer rigidity and unconditional majority}
\label{sec:integer}

\begin{corollary}[Integer profile and exact right-hand plateau]
\label{cor:integer}
Assume Conjecture~\CGref{conj:affine_chowla}. For fixed $m\ge2$,
\begin{equation}
\frac1X C_X\!\left(\log m+\frac cX\right)
\longrightarrow
\begin{cases}
0,&c<0,\\
\lambda(m),&0\le c\le1/m,\\
\lambda(m)/(mc),&c>1/m.
\end{cases}
\label{eq:o19}
\end{equation}
Moreover, for every fixed $0<c<1/m$,
\begin{equation}
C_X\!\left(\log m+\frac cX\right)=\lambda(m)X
\label{eq:o20}
\end{equation}
for all sufficiently large $X$.
\end{corollary}

\begin{proof}
The limit is Theorem~\CGref{thm:profile} with $q=1$. For the exact statement,
$mX(e^{c/X}-1)\to mc<1$, so eventually no correction occurs for any
$n\le X$.
\end{proof}

\begin{theorem}[Unconditional right-hand natural majority window]
\label{thm:naturalmajority}
Let $m\ge1$ and $0<c<2/m$. Then
\begin{equation}
\liminf_{X\to\infty}
\lambda(m)\frac1X
C_X\!\left(\log m+\frac cX\right)
\ge
2\min\!\left(1,\frac1{mc}\right)-1>0.
\label{eq:naturalmajority}
\end{equation}
\end{theorem}

\begin{proof}
Let
\begin{equation}
I_{0,X}:=\{1\le n\le X:m(e^{c/X}-1)n<1\}.
\label{eq:o21}
\end{equation}
Every summand on $I_{0,X}$ equals $\lambda(m)$, while every remaining term
has modulus at most one. Therefore
\begin{equation}
\lambda(m)C_X\!\left(\log m+\frac cX\right)\ge2|I_{0,X}|-X.
\label{eq:o22}
\end{equation}
Since
\begin{equation}
\frac{|I_{0,X}|}{X}\to\min\!\left(1,\frac1{mc}\right),
\label{eq:o23}
\end{equation}
the result follows.
\end{proof}

\begin{remark}[Signed natural majority bound]
Theorem~\CGref{thm:naturalmajority} is a signed lower bound. In particular,
absolute non-cancellation follows immediately.
\end{remark}

\begin{theorem}[Unconditional logarithmic mesoscopic majority]
\label{thm:logmajority}
Let $(p,q)=1$ and let $c_X>0$ satisfy
\begin{equation}
\frac{\log c_X}{\log X}\longrightarrow\theta,
\qquad 0\le\theta<1.
\label{eq:o24}
\end{equation}
Then
\begin{equation}
\liminf_{X\to\infty}
\lambda(pq)\,
C_X^{\log}\!\left(\log\frac pq+\frac{c_X}{X}\right)
\ge
\frac{2(1-\theta)}q-1.
\label{eq:logmajority}
\end{equation}
In particular, when $q=1$ and $\theta<1/2$, writing $m=p$,
\begin{equation}
\liminf_{X\to\infty}
\lambda(m)C_X^{\log}\!\left(\log m+\frac{c_X}{X}\right)
\ge1-2\theta>0.
\label{eq:o25}
\end{equation}
\end{theorem}

\begin{proof}
On the proportional class $n=qk$, the zero-correction condition is
\begin{equation}
p(e^{c_X/X}-1)k<1.
\label{eq:o26}
\end{equation}
Let $K_X$ be the largest admissible $k$, truncated by $k\le X/q$. Since
$c_X=X^{\theta+o(1)}$ and $\theta<1$,
\begin{equation}
K_X=X^{1-\theta+o(1)}.
\label{eq:o27}
\end{equation}
The normalized harmonic mass of this block is
\begin{equation}
B_X=
\frac1{H_X}\sum_{k\le K_X}\frac1{qk}
=
\frac{1-\theta}{q}+o(1).
\label{eq:o28}
\end{equation}
After multiplication by $\lambda(pq)$, all terms on this block are positive.
The complement has total normalized harmonic mass at most $1-B_X$, so
\begin{equation}
\lambda(pq)C_X^{\log}\ge B_X-(1-B_X)=2B_X-1.
\label{eq:o29}
\end{equation}
This proves \CGeqref{eq:logmajority}.
\end{proof}

\begin{remark}[Threshold of the elementary logarithmic majority mechanism]
For $q=1$, the lower bound in Theorem~\CGref{thm:logmajority} is positive
exactly when $\theta<1/2$. Thus $1/2$ is the threshold exponent for this
elementary majority mechanism, not a claimed transition exponent for the
correlation itself. More precisely, every sequence
$c_X=X^{\theta+o(1)}$ with $0\le\theta<1/2$ is covered, at the displacement
$c_X/X=X^{-1+\theta+o(1)}$; in particular this includes
$X^{-1/2-\varepsilon+o(1)}$ for every $\varepsilon>0$.
For $q\ge2$, this elementary majority bound yields no positive lower bound.
\end{remark}

\begin{remark}[Positive fixed logarithmic branch from the majority theorem]
For $q=1$ and $\theta=0$, Theorem~\CGref{thm:logmajority} gives
$\liminf\lambda(m)C_X^{\log}\ge1$, while $|C_X^{\log}|\le1$. Hence
\begin{equation}
C_X^{\log}\!\left(\log m+\frac{c_X}{X}\right)\to\lambda(m)
\label{eq:o30}
\end{equation}
for every positive $c_X$ with $\log c_X/\log X\to0$. No logarithmic Chowla
input is needed for this integer positive branch.
\end{remark}

\begin{remark}[The endpoint $c=1/m$ is not an exact finite plateau]
For the integer right-hand layer,
\begin{equation}
Xm(e^{1/(mX)}-1)
=
1+\frac1{2mX}+O(X^{-2})>1.
\label{eq:o31}
\end{equation}
Thus $c=1/m$ belongs to the closed limiting plateau. Call $c$ in the
\emph{exact finite plateau} if $C_X(\log m+c/X)=\lambda(m)X$ for all sufficiently
large $X$. At $c=1/m$ only the term $n=X$ changes, and
$C_X(\log m+\tfrac1{mX})=\lambda(m)(X-1)+\lambda(X)\lambda(mX+1)$; this differs from
$\lambda(m)X$ whenever $\lambda(X)\lambda(mX+1)=-\lambda(m)$, which happens for
infinitely many $X$ because $\sum_{X\le Y}\lambda(X)\lambda(mX+1)=o(Y)$
(Theorem~\CGref{thm:oai}). Hence $c=1/m$ is not in the exact finite plateau. For $m=2$, $X=10^6$, the regenerated computation gives
$C_X=-999998$.
\end{remark}

\begin{proposition}[Exact left-hand $c$-independence interval]
\label{prop:left}
Let $m\ge2$ and fix $-1/m<c<0$. For all sufficiently large $X$,
\begin{equation}
C_X\!\left(\log m+\frac cX\right)
=
\sum_{n\le X}\lambda(n)\lambda(mn-1).
\label{eq:o32}
\end{equation}
This states exact independence with respect to $c$ on that finite interval;
no asymptotic value of the right-hand sum is asserted.
\end{proposition}

\begin{proof}
Put $\beta_X=m(1-e^{c/X})$. Since $X\beta_X\to m|c|<1$, eventually
$0<\beta_Xn<1$ for all $n\le X$. Hence
$\lceil\beta_Xn\rceil=1$ and $\lfloor me^{c/X}n\rfloor=mn-1$.
\end{proof}

\begin{proposition}[Locally finite step structure]
Fix $m\ge2$ and $X\in\mathbb N$. The map
\begin{equation}
c\longmapsto C_X\!\left(\log m+\frac cX\right)
\label{eq:o33}
\end{equation}
is locally constant except at explicit crossing values. Every compact
$c$-interval contains only finitely many jumps. For sufficiently negative
$c$, disappearance points where $\lfloor me^{c/X}n\rfloor<1$ are included
among these crossings.
\end{proposition}

\begin{proof}
For each fixed $n\le X$, the integer $\lfloor me^{c/X}n\rfloor$ changes only
when the continuous function $me^{c/X}n$ crosses an integer. On a compact
$c$-interval its range is bounded, hence only finitely many crossings occur.
There are only finitely many $n\le X$.
\end{proof}

\section{Unconditional natural profile}
\label{sec:unconditional}

Conjecture~\CGref{conj:affine_chowla} is used in Sections~\CGref{sec:natural} and~\CGref{sec:integer} only for finitely many fixed pairs of non-proportional affine forms, namely those produced by the residue-class blocks of Proposition~\CGref{prop:decomposition} and of the proof of Theorem~\CGref{thm:profile}. The following theorem supplies exactly this input, with a rate.

\begin{theorem}[Ordinary two-point affine Chowla {\cite[Theorem~1.1, p.~3]{OpenAI2026}}]
\label{thm:oai}
There is an absolute constant $\kappa>0$ with the following property. Let $a_1,a_2\ge1$ and $b_1,b_2\ge0$ be fixed integers with $a_1b_2-a_2b_1\ne0$. Then there is $C=C_{a_1,a_2,b_1,b_2}$ such that
\begin{equation}
\Bigl|\sum_{n\le X}\lambda(a_1n+b_1)\lambda(a_2n+b_2)\Bigr|\le C\,\frac{X}{(\log X)^{\kappa}}\qquad(X\ge3,\ X\in\mathbb R).
\label{eq:oai}
\end{equation}
\end{theorem}

The constants in \CGeqref{eq:oai} may be ineffective, and no uniformity in coefficients growing with $X$ is asserted in \cite[p.~3]{OpenAI2026}; none is used below.

\begin{lemma}[Macroscopic blocks with a rate]
\label{lem:macro_rate}
Let $\alpha_1,\alpha_2\ge1$ and $\beta_1,\beta_2\in\mathbb Z$ be fixed with $\alpha_1\beta_2-\alpha_2\beta_1\ne0$. Let $0\le A\le B<\infty$ be fixed and $A_X=AX+O(1)$, $B_X=BX+O(1)$. Then
\begin{equation}
\sum_{\substack{A_X<k\le B_X\\ \alpha_1k+\beta_1\ge1,\ \alpha_2k+\beta_2\ge1}}\lambda(\alpha_1k+\beta_1)\lambda(\alpha_2k+\beta_2)=O\Bigl(\frac{X}{(\log X)^{\kappa}}\Bigr).
\label{eq:macro_rate}
\end{equation}
\end{lemma}

\begin{proof}
Choose $k_0\ge0$ with $\alpha_ik_0+\beta_i\ge0$ for $i=1,2$ and substitute $k=\ell+k_0$. The forms become $\alpha_i\ell+(\alpha_ik_0+\beta_i)$, with non-negative intercepts and the same determinant $\alpha_1\beta_2-\alpha_2\beta_1\ne0$. The positivity conditions remove only $O(1)$ terms. The sum is a difference of two prefix sums at the cutoffs $A_X$ and $B_X$. Theorem~\CGref{thm:oai} applies to each cutoff $\asymp X$; when $A=0$ the lower prefix sum has $A_X=O(1)$ terms and is $O(1)$.
\end{proof}

\begin{theorem}[Rational natural microscopic profile, unconditional]
\label{thm:profile_uncond}
Let $(p,q)=1$ and fix $c\in\mathbb R$. Then
\begin{equation}
\frac1X\,C_X\!\Bigl(\log\frac pq+\frac cX\Bigr)=\Psi_{p,q}(c)+O_{p,q,c}\bigl((\log X)^{-\kappa}\bigr),
\label{eq:profile_uncond}
\end{equation}
with $\Psi_{p,q}$ as in \CGeqref{eq:profile} and $\kappa$ the absolute constant of Theorem~\CGref{thm:oai}.
\end{theorem}

\begin{proof}
For $c=0$ use Proposition~\CGref{prop:decomposition}: the proportional class contributes $\lambda(pq)\lfloor X/q\rfloor$ and each of the $q-1$ remaining classes is a sum of the shape \CGeqref{eq:macro_rate} with $A=0$, $B=1/q$. Let $c\ne0$. In the proof of Theorem~\CGref{thm:profile} the number of pairs $(a,r)$ is bounded in terms of $p,q,c$. Since $Xp(e^{c/X}-1)=pc+O(X^{-1})$, the endpoints in \CGeqref{eq:blockinterval} are $X\,\frac{r-\{pa/q\}}{pc}+O(1)$; after intersection with $0\le k\le X/q$ every block is an interval $(A_X,B_X]$ with $A_X=AX+O(1)$, $B_X=BX+O(1)$ for fixed $A\le B$. Every block other than $(a,r)=(0,0)$ consists of two non-proportional fixed affine forms, so Lemma~\CGref{lem:macro_rate} bounds it by $O(X(\log X)^{-\kappa})$. The proportional block has cardinality $R_X(c)=X\min(1/q,1/(pc))\mathbf 1_{c>0}+O(1)$, and each of its terms equals $\lambda(pq)$.
\end{proof}

\begin{corollary}[Unconditional forms of Theorem~\CGref{thm:fixed}, Corollary~\CGref{cor:integer} and Proposition~\CGref{prop:left}]
\label{cor:uncond}
Let $(p,q)=1$ and $m\ge2$.
\begin{enumerate}
\item $C_X(\log(p/q))=\lambda(pq)X/q+O(X(\log X)^{-\kappa})$.
\item The limit in Corollary~\CGref{cor:integer} holds unconditionally, with error term $O_{m,c}((\log X)^{-\kappa})$.
\item For fixed $-1/m<c<0$ and all sufficiently large $X$,
\begin{equation}
C_X\!\Bigl(\log m+\frac cX\Bigr)=\sum_{n\le X}\lambda(n)\lambda(mn-1)=O\Bigl(\frac X{(\log X)^{\kappa}}\Bigr).
\label{eq:left_uncond}
\end{equation}
\end{enumerate}
\end{corollary}

\begin{proof}
Items (1) and (2) are the cases $c=0$ and $q=1$ of Theorem~\CGref{thm:profile_uncond}. For (3), the identity is Proposition~\CGref{prop:left}, and the forms $n$, $mn-1$ have determinant $-1\ne0$, so Lemma~\CGref{lem:macro_rate} applies.
\end{proof}

\begin{remark}[The trivial resonance is binary Chowla]
\label{rem:trivial_resonance}
Take $p=q=1$ and fix $-1<c<0$. Since $X(1-e^{c/X})\to|c|<1$, for $X\ge X_0(c)$ one has $0<n(1-e^{c/X})<1$ for every $n\le X$, hence $\lfloor e^{c/X}n\rfloor=n-1$, and the term $n=1$ is omitted by definition. Thus, for every $X\ge X_0(c)$,
\begin{equation}
C_X\!\Bigl(\frac cX\Bigr)=\sum_{2\le n\le X}\lambda(n-1)\lambda(n),
\label{eq:trivial_resonance}
\end{equation}
while $\Psi_{1,1}(c)=0$. The profile of Theorem~\CGref{thm:profile} at the single point $(p,q,c)=(1,1,c)$ is therefore, for every large $X$ and not only in the limit, the shift-one Chowla sum. No argument that avoids two-point input can produce Theorem~\CGref{thm:profile}; the residue-class decomposition isolates the main term exactly, and the remainder is a finite set of fixed binary correlations, which Theorem~\CGref{thm:oai} controls.
\end{remark}


\section{The \texorpdfstring{$k$}{k}-point resonance family}
\label{sec:kpoint}

\subsection{Definition and exact decomposition}
\label{subsec:kpoint_def}

For $k\ge2$ and $u=(u_1,\dots,u_{k-1})\in\mathbb R^{k-1}$ put
\begin{equation}
\begin{gathered}
C^{(k)}_X(u):=\sum_{\substack{1\le n\le X\\ \lfloor e^{u_i}n\rfloor\ge1\ \forall i}}\lambda(n)\prod_{i=1}^{k-1}\lambda\bigl(\lfloor e^{u_i}n\rfloor\bigr),
\\
C^{(k),\log}_X(u):=\frac1{H_X}\sum_{\substack{1\le n\le X\\ \lfloor e^{u_i}n\rfloor\ge1\ \forall i}}\frac{\lambda(n)\prod_{i}\lambda(\lfloor e^{u_i}n\rfloor)}{n}.
\end{gathered}
\label{eq:Ck}
\end{equation}
For $k=2$ these are $C_X(u)$ and $C^{\log}_X(u)$ of Section~\CGref{sec:input}. Fix coprime $p_i,q_i\ge1$ and $c_i\in\mathbb R$ and set
\begin{equation}
u_i=\log\frac{p_i}{q_i}+\frac{c_i}{X},\qquad L:=\operatorname{lcm}(q_1,\dots,q_{k-1}),\qquad \delta_i:=\frac{p_i}{q_i}\bigl(e^{c_i/X}-1\bigr).
\label{eq:uk}
\end{equation}
Write $n=Lm+a$ with $0\le a<L$, $m\ge0$, and
\begin{equation}
\theta_{a,i}:=\Bigl\{\frac{p_ia}{q_i}\Bigr\},\qquad b_{a,i}:=\Bigl\lfloor\frac{p_ia}{q_i}\Bigr\rfloor,\qquad A_i:=\frac{p_iL}{q_i}\in\mathbb N .
\label{eq:theta_k}
\end{equation}

\begin{lemma}[Floors on residue classes]
\label{lem:floor_k}
For $n=Lm+a\le X$,
\begin{equation}
\lfloor e^{u_i}n\rfloor=A_im+b_{a,i}+r_i(n),\qquad r_i(n):=\bigl\lfloor\theta_{a,i}+\delta_in\bigr\rfloor,
\label{eq:floor_k}
\end{equation}
and $|r_i(n)|\le\frac{p_i}{q_i}|c_i|e^{|c_i|}+1$. In particular $r_i(n)$ ranges over a finite set $R_i$ depending only on $(p_i,q_i,c_i)$.
\end{lemma}

\begin{proof}
$e^{u_i}n=\frac{p_i}{q_i}(Lm+a)+\delta_in=A_im+b_{a,i}+\theta_{a,i}+\delta_in$, and $|\delta_in|\le\frac{p_i}{q_i}|e^{c_i/X}-1|X\le\frac{p_i}{q_i}|c_i|e^{|c_i|}$ for $n\le X$.
\end{proof}

For $0\le a<L$ and $r=(r_1,\dots,r_{k-1})\in R:=\prod_iR_i$ let $I_{a,r}(X)$ be the set of $m\ge0$ with $1\le Lm+a\le X$, $f_{a,r,i}(m)\ge1$ for all $i$, and $r_i(Lm+a)=r_i$ for all $i$, where
\begin{equation}
f_{a,r,0}(m):=Lm+a,\qquad f_{a,r,i}(m):=A_im+b_{a,i}+r_i\quad(1\le i\le k-1).
\label{eq:forms_k}
\end{equation}
Since $n\mapsto\theta_{a,i}+\delta_in$ is affine, $I_{a,r}(X)$ is an interval of integers, and by Lemma~\CGref{lem:floor_k}
\begin{equation}
C^{(k)}_X(u)=\sum_{a=0}^{L-1}\sum_{r\in R}\ \sum_{m\in I_{a,r}(X)}\ \prod_{i=0}^{k-1}\lambda\bigl(f_{a,r,i}(m)\bigr)
\label{eq:dec0_k}
\end{equation}
exactly, for every $X$.

\begin{definition}[Proportionality classes]
\label{def:classes}
For fixed $(a,r)$ write $f_{a,r,i}(m)=\mathcal A_im+\mathcal B_i$, let $d_i:=\gcd(\mathcal A_i,|\mathcal B_i|)$ (with $\gcd(\mathcal A,0)=\mathcal A$), and let $(\alpha_i,\beta_i):=(\mathcal A_i,\mathcal B_i)/d_i$, a primitive vector with $\alpha_i\ge1$. Declare $i\sim j$ if $(\alpha_i,\beta_i)=(\alpha_j,\beta_j)$, equivalently $\mathcal A_i\mathcal B_j=\mathcal A_j\mathcal B_i$. For a class $\mathcal C$ let $g_{\mathcal C}(m):=\alpha_{\mathcal C}m+\beta_{\mathcal C}$ be its primitive form, let $\mathrm{Odd}(a,r)$ be the set of classes of odd cardinality and $k'(a,r):=|\mathrm{Odd}(a,r)|$. The block is \emph{resonant} if $k'(a,r)=0$. Put $\varepsilon_{a,r}:=\lambda(d_0d_1\cdots d_{k-1})\in\{\pm1\}$.
\end{definition}

\begin{theorem}[Exact decomposition of the $k$-point family]
\label{thm:dec_k}
For every $X\ge1$ and every choice \CGeqref{eq:uk},
\begin{equation}
C^{(k)}_X(u)=\sum_{\substack{(a,r)\\ k'(a,r)=0}}\varepsilon_{a,r}\,\bigl|I_{a,r}(X)\bigr|
+\sum_{\substack{(a,r)\\ k'(a,r)\ge1}}\varepsilon_{a,r}\sum_{m\in I_{a,r}(X)}\ \prod_{\mathcal C\in\mathrm{Odd}(a,r)}\lambda\bigl(g_{\mathcal C}(m)\bigr).
\label{eq:dec_k}
\end{equation}
The forms $g_{\mathcal C}$, $\mathcal C\in\mathrm{Odd}(a,r)$, are primitive, pairwise non-proportional and positive on $I_{a,r}(X)$, and $k'(a,r)\equiv k\pmod 2$.
\end{theorem}

\begin{proof}
On $I_{a,r}(X)$ every $f_{a,r,i}(m)\ge1$, hence $g_{\mathcal C}(m)\ge1$ and $\lambda(f_{a,r,i}(m))=\lambda(d_i)\lambda(g_{\mathcal C}(m))$ by complete multiplicativity. Since $\lambda^2=1$,
\begin{equation}
\prod_{i=0}^{k-1}\lambda(f_{a,r,i}(m))=\lambda\Bigl(\prod_i d_i\Bigr)\prod_{\mathcal C}\lambda\bigl(g_{\mathcal C}(m)\bigr)^{|\mathcal C|}=\varepsilon_{a,r}\prod_{\mathcal C\in\mathrm{Odd}(a,r)}\lambda\bigl(g_{\mathcal C}(m)\bigr).
\label{eq:collapse}
\end{equation}
Distinct classes have distinct primitive vectors and are therefore non-proportional, and $\sum_{\mathcal C}|\mathcal C|=k$ gives the parity. Substitute into \CGeqref{eq:dec0_k}.
\end{proof}

\begin{lemma}[Block lengths]
\label{lem:len_k}
Let $J_{a,r}(c)\subseteq[0,1]$ be the set of $t$ with $r_i\le\theta_{a,i}+\frac{p_ic_i}{q_i}t<r_i+1$ for all $i$. It is an interval, possibly empty, and $|I_{a,r}(X)|=\frac XL|J_{a,r}(c)|+O_{p,q,c}(1)$. More precisely, $\{m:\ m\in I_{a,r}(X)\}=(A_X,B_X]\cap\mathbb Z$ with $A_X=AX+O(1)$, $B_X=BX+O(1)$ for fixed $A\le B$.
\end{lemma}

\begin{proof}
For $n\le X$, $\delta_in=\frac{p_ic_i}{q_i}\frac nX(1+O(|c_i|/X))$, so the endpoints of $\{Lm/X: m\in I_{a,r}(X)\}$ differ from those of $J_{a,r}(c)$ by $O(1/X)$, and the conditions $f_{a,r,i}(m)\ge1$ remove $O(1)$ values of $m$.
\end{proof}

\begin{definition}[$k$-point profile]
\label{def:profile_k}
\begin{equation}
\Psi^{(k)}_{p,q}(c):=\frac1L\sum_{\substack{(a,r)\\ k'(a,r)=0}}\varepsilon_{a,r}\,\bigl|J_{a,r}(c)\bigr| .
\label{eq:profile_k}
\end{equation}
For odd $k$ there are no resonant blocks, since $k'\equiv k\pmod2$, and $\Psi^{(k)}\equiv0$.
\end{definition}

\begin{remark}[$k=2$]
\label{rem:k2}
For $k=2$, $L=q$, and $0\sim1$ holds if and only if $q\mid a$ and $r=0$, that is $(a,r)=(0,0)$, with $d_0=q$, $d_1=p$, $\varepsilon_{0,0}=\lambda(pq)$, and $|J_{0,0}(c)|=\min(1,q/(pc))$ for $c\ge0$, $|J_{0,0}(c)|=0$ for $c<0$. Definition~\CGref{def:profile_k} recovers $\Psi_{p,q}$ of Theorem~\CGref{thm:profile}, and Theorem~\CGref{thm:dec_k} recovers the block structure of its proof.
\end{remark}

\subsection{What each order of Chowla closes}
\label{subsec:kpoint_suff}

Say that $\mathrm{Ch}_{k'}$ holds if, for every tuple of $k'$ pairwise non-proportional affine forms $\alpha_jm+\beta_j$ ($\alpha_j\ge1$, $\beta_j\in\mathbb Z$), $\sum_{m\le Y}\prod_{j}\lambda(\alpha_jm+\beta_j)=o(Y)$, the finitely many $m$ with some form $\le0$ being omitted.

\begin{theorem}[Sufficiency]
\label{thm:suf_k}
In \CGeqref{eq:dec_k}:
\begin{enumerate}
\item the pieces with $k'(a,r)=1$ are $O_{p,q,c}(X\exp(-c_0\sqrt{\log X}))$ for an absolute $c_0>0$;
\item the pieces with $k'(a,r)=2$ are $O_{p,q,c}(X(\log X)^{-\kappa})$;
\item if $\mathrm{Ch}_{k'}$ holds for every $3\le k'\le k$ with $k'\equiv k\pmod2$, then
\begin{equation}
\frac1XC^{(k)}_X(u)\longrightarrow\Psi^{(k)}_{p,q}(c)\qquad(X\to\infty).
\label{eq:limit_k}
\end{equation}
\end{enumerate}
\end{theorem}

\begin{proof}
By Lemma~\CGref{lem:len_k} each piece is $\varepsilon_{a,r}$ times a sum over $(A_X,B_X]$, a difference of two prefix sums. (1) For one primitive form $\alpha m+\beta$ this is a sum of $\lambda$ over a fixed progression, bounded by the prime number theorem in arithmetic progressions of fixed modulus. (2) Two non-proportional primitive forms: Lemma~\CGref{lem:macro_rate}. (3) The pieces with $k'\ge3$ are $o(X)$ by $\mathrm{Ch}_{k'}$, and Lemma~\CGref{lem:len_k} gives the resonant part.
\end{proof}

\begin{proposition}[Exact criterion for unconditional closure]
\label{prop:criterion}
Fix $(p,q,c)$. If every block $(a,r)$ with $|J_{a,r}(c)|>0$ has $k'(a,r)\le2$, then \CGeqref{eq:limit_k} holds unconditionally, with error $O((\log X)^{-\kappa})$. Conversely, if some block with $|J_{a,r}(c)|>0$ has $k'(a,r)\ge3$, then Theorem~\CGref{thm:suf_k}(1),(2) do not cover it, and \CGeqref{eq:limit_k} at $(p,q,c)$ is equivalent to the statement that the sum over all such blocks is $o(X)$; it does not follow that each block separately satisfies its $k'$-point estimate.
\end{proposition}

\begin{proof}
Blocks with $|J_{a,r}(c)|=0$ have $O(1)$ elements. The others are covered by Theorem~\CGref{thm:suf_k}(1),(2) exactly when $k'\le2$.
\end{proof}

\begin{example}[$k=4$, closed unconditionally]
\label{ex:k4}
Take $(p_i/q_i)=(2,3,6)$ and $c=(1,c_2,1)$ with $0<c_2<\tfrac13$. Then $L=1$, $a=0$, $r_2=0$ on $[0,1]$, $r_1=\lfloor2t\rfloor$ and $r_3=\lfloor6t\rfloor$. The blocks in $t=n/X$, with the proportionality classes of the forms $(m,\,2m+r_1,\,3m,\,6m+r_3)$ indexed $0,1,2,3$, are:
\begin{center}
\small
\begin{tabular}{llll}
\hline
$t$ & $r$ & classes & piece\\
\hline
$[0,\tfrac16)$ & $(0,0,0)$ & $\{0,1,2,3\}$ & resonant, $\varepsilon=\lambda(36)=1$\\
$[\tfrac16,\tfrac13)$ & $(0,0,1)$ & $\{0,1,2\},\{3\}$ & $\lambda(6)\lambda(m)\lambda(6m+1)$, $k'=2$\\
$[\tfrac13,\tfrac12)$ & $(0,0,2)$ & $\{0,1,2\},\{3\}$ & $\lambda(12)\lambda(m)\lambda(3m+1)$, $k'=2$\\
$[\tfrac12,\tfrac23)$ & $(1,0,3)$ & $\{0,2\},\{1,3\}$ & resonant, $\varepsilon=\lambda(9)=1$\\
$[\tfrac23,\tfrac56)$ & $(1,0,4)$ & $\{0,2\},\{1\},\{3\}$ & $\lambda(6)\lambda(2m+1)\lambda(3m+2)$, $k'=2$\\
$[\tfrac56,1)$ & $(1,0,5)$ & $\{0,2\},\{1\},\{3\}$ & $\lambda(3)\lambda(2m+1)\lambda(6m+5)$, $k'=2$\\
\hline
\end{tabular}
\end{center}
By Proposition~\CGref{prop:criterion}, unconditionally,
\begin{equation}
\frac1X\,C^{(4)}_X\Bigl(\log2+\frac1X,\ \log3+\frac{c_2}X,\ \log6+\frac1X\Bigr)=\frac13+O\bigl((\log X)^{-\kappa}\bigr).
\label{eq:ex4}
\end{equation}
The second plateau comes from the block $(1,0,3)$, where the four forms pair into $(m,3m)$ and $(2m+1,6m+3)$ without all four being proportional; this pairing has no analogue for $k=2$.
\end{example}

\subsection{Plateaus, majority and the left-hand identity at integer resonances}
\label{subsec:kpoint_integer}

\begin{theorem}[Integer resonances at level $k$]
\label{thm:integer_k}
Let $m_1,\dots,m_{k-1}\ge1$ be integers and $u_i=\log m_i+c_i/X$.
\begin{enumerate}
\item If $0<c_i<1/m_i$ for all $i$, then for $X\ge X_0(m,c)$
\begin{equation}
C^{(k)}_X(u)=\lambda(m_1\cdots m_{k-1})\sum_{n\le X}\lambda(n)^k=
\begin{cases}\lambda(m_1\cdots m_{k-1})\lfloor X\rfloor,& k\ \text{even},\\[1mm]
O\bigl(Xe^{-c_0\sqrt{\log X}}\bigr),& k\ \text{odd}.\end{cases}
\label{eq:plateau_k}
\end{equation}
\item If $k$ is even and $c_i>0$ for all $i$, then
\begin{equation}
\liminf_{X\to\infty}\ \lambda(m_1\cdots m_{k-1})\frac{C^{(k)}_X(u)}{X}\ \ge\ 2\min\Bigl(1,\frac1{\max_im_ic_i}\Bigr)-1 .
\label{eq:majority_k}
\end{equation}
\item If $-1/m_i<c_i<0$ for all $i$, then for $X\ge X_0(m,c)$
\begin{equation}
C^{(k)}_X(u)=\sum_{\substack{n\le X\\ m_in-1\ge1\ \forall i}}\lambda(n)\prod_{i=1}^{k-1}\lambda(m_in-1).
\label{eq:left_k}
\end{equation}
\end{enumerate}
No Chowla input is used.
\end{theorem}

\begin{proof}
$\lfloor m_ie^{c_i/X}n\rfloor=m_in+\lfloor m_i(e^{c_i/X}-1)n\rfloor$ for $c_i>0$. (1) Since $m_iX(e^{c_i/X}-1)\to m_ic_i<1$, eventually the correction vanishes for every $n\le X$, and $\lambda(n)\prod_i\lambda(m_in)=\lambda(\prod m_i)\lambda(n)^k$; for odd $k$ use the prime number theorem with classical error term. (2) On $n<\min_i\{m_i(e^{c_i/X}-1)\}^{-1}$ every correction vanishes and every summand equals $\lambda(\prod m_i)\lambda(n)^k=\lambda(\prod m_i)$; the proportion of such $n\le X$ tends to $\min(1,1/\max_im_ic_i)$, and every other summand has modulus one. (3) Since $m_iX(1-e^{c_i/X})\to m_i|c_i|<1$, eventually $0<m_i(1-e^{c_i/X})n<1$ for $n\le X$, so $\lfloor m_ie^{c_i/X}n\rfloor=m_in-1$.
\end{proof}

Parts (1) and (2) extend Corollary~\CGref{cor:integer} and Theorem~\CGref{thm:naturalmajority} to every $k$; part (3) extends Proposition~\CGref{prop:left}. If the $m_i$ are pairwise distinct, the $k$ forms $n,m_1n-1,\dots,m_{k-1}n-1$ in \CGeqref{eq:left_k} are pairwise non-proportional, so $k'=k$ and $\Psi^{(k)}=0$ there. The left-hand plateau is thus, for every large $X$, a single $k$-point Chowla sum.

\begin{corollary}[Level three at $(2,3)$ is ordinary three-point Chowla]
\label{cor:level3_23}
For every $c=(c_1,c_2)\in(-\tfrac12,0)\times(-\tfrac13,0)$ and every $X\ge X_0(c)$,
\begin{equation}
\begin{aligned}
C^{(3)}_X\Bigl(\log2+\frac{c_1}X,\ \log3+\frac{c_2}X\Bigr)&=\sum_{n\le X}\lambda(n)\lambda(2n-1)\lambda(3n-1)\\
&=\sum_{\substack{3\le M\le 6X-3\\ M\equiv3\ (\mathrm{mod}\ 6)}}\lambda(M)\lambda(M+1)\lambda(M+3),
\end{aligned}
\label{eq:level3_23}
\end{equation}
while $\Psi^{(3)}\equiv0$. Consequently the level-three profile on this open set of $c$, for all large $X$, is equivalent to
\begin{equation}
\sum_{\substack{M\le Y\\ M\equiv3\ (\mathrm{mod}\ 6)}}\lambda(M)\lambda(M+1)\lambda(M+3)=o(Y).
\label{eq:chowla013}
\end{equation}
\end{corollary}

\begin{proof}
The first identity is Theorem~\CGref{thm:integer_k}(3). Put $N=6n$. Then $3(2n-1)=N-3$ and $2(3n-1)=N-2$, so $\lambda(n)=\lambda(6)\lambda(N)=\lambda(N)$, $\lambda(2n-1)=\lambda(3)\lambda(N-3)$ and $\lambda(3n-1)=\lambda(2)\lambda(N-2)$; since $\lambda(2)\lambda(3)=1$ the product is $\lambda(N)\lambda(N-2)\lambda(N-3)$. Put $M=N-3$.
\end{proof}

The logarithmic analogue of \CGeqref{eq:chowla013} holds by \cite[Theorem~1.1, p.~2]{TaoTeravainen2018}; the ordinary average \CGeqref{eq:chowla013} at every scale is not known, and \cite{OpenAI2026} treats only two points \cite[pp.~3--4]{OpenAI2026}.

\subsection{Necessity}
\label{subsec:kpoint_nec}

\begin{theorem}[Necessity]
\label{thm:nec_k}
Fix $k\ge2$ and $a_1,\dots,a_{k-1}\ge1$. For $r\in\mathbb Z_{\ge0}^{k-1}$ put
\begin{equation}
P_r(Y):=\sum_{m\le Y}\lambda(m)\prod_{i=1}^{k-1}\lambda(a_im+r_i),\qquad \pi_r:=\varepsilon_{0,r}\,\mathbf 1\bigl[k'(0,r)=0\bigr],
\label{eq:Pr}
\end{equation}
with the classes of Definition~\CGref{def:classes} for $(p_i,q_i)=(a_i,1)$. Suppose that \CGeqref{eq:limit_k} holds for $(p_i,q_i)=(a_i,1)$ and every $c\in(0,\infty)^{k-1}$. Then $P_r(Y)=\pi_rY+o(Y)$ for every $r\in\mathbb Z_{\ge0}^{k-1}$. In particular, if the forms $m,a_1m+r_1,\dots,a_{k-1}m+r_{k-1}$ are pairwise non-proportional, then $P_r(Y)=o(Y)$.
\end{theorem}

\begin{proof}
Induction on $|r|:=\sum_ir_i$. For $r=0$, $P_0(Y)=\lambda(\prod a_i)\sum_{m\le Y}\lambda(m)^k$, which is $\pi_0Y+o(Y)$ by the prime number theorem when $k$ is odd and trivially when $k$ is even. Let $r\ne0$ and choose $c_i:=(r_i+1)/a_i$. Then $\delta_im=(r_i+1)\frac mX(1+\eta_i)$ with $0<\eta_i\le c_i/X$, so apart from $O(1)$ values of $m$ near $X$ every nonempty block $I_{r'}(X)$ has $r'\le r$ componentwise, and
\begin{equation}
I_r(X)=\bigl[\alpha_rX,\ X\bigr]+O(1),\qquad \alpha_r:=\max_i\frac{r_i}{r_i+1}\in\bigl[\tfrac12,1\bigr).
\label{eq:last_block}
\end{equation}
Every other nonempty block has $|r'|<|r|$, is an interval $(\alpha'X,\beta'X]+O(1)$, and by induction contributes $\pi_{r'}|J_{r'}(c)|X+o(X)$. The hypothesis \CGeqref{eq:limit_k} and \CGeqref{eq:profile_k} leave $\sum_{m\in I_r(X)}\lambda(m)\prod_i\lambda(a_im+r_i)=\pi_r(1-\alpha_r)X+o(X)$, that is
\begin{equation}
P_r(Y)-P_r(\alpha_rY)=\pi_r(1-\alpha_r)Y+o(Y)\qquad(Y\to\infty).
\label{eq:suffix}
\end{equation}
Given $\epsilon>0$ take $Y_0$ with error $\le\epsilon Y$ in \CGeqref{eq:suffix} for $Y\ge Y_0$, and $J$ with $\alpha_r^{J+1}Y<Y_0\le\alpha_r^JY$. Telescoping gives $|P_r(Y)-\pi_r(1-\alpha_r^{J+1})Y|\le\epsilon Y/(1-\alpha_r)+Y_0$, hence $\limsup_Y|P_r(Y)-\pi_rY|/Y\le\epsilon/(1-\alpha_r)$. If the forms are pairwise non-proportional, all classes are singletons and $\pi_r=0$.
\end{proof}

For instance, the level-three profile at $(2,3)$ on $(0,\infty)^2$ implies
\begin{equation}
\sum_{m\le X}\lambda(m)\lambda(2m+1)\lambda(3m+1)=o(X):
\label{eq:231}
\end{equation}
with $c=(1,\tfrac23)$ the blocks are $r=(0,0)$ on the range $t\in[0,\tfrac12)$, where $\lambda(m)\lambda(2m)\lambda(3m)=\lambda(m)$, and $r=(1,1)$ on $[\tfrac12,1]$. Together with Theorem~\CGref{thm:suf_k} and Corollary~\CGref{cor:level3_23}, for $q_i=1$ and $c\in(0,\infty)^{k-1}$ the level-$k$ family is equivalent to $k$-point Chowla on the tuples it contains (Theorem~\CGref{thm:nec_k}); in general, for $k\ge3$, it is an \emph{exact} reformulation and not a reduction to two points.

\begin{theorem}[Odd levels, logarithmic average]
\label{thm:log_k}
Let $k\ge3$ be odd. For every rational resonance \CGeqref{eq:uk} and every $c\in\mathbb R^{k-1}$, $C^{(k),\log}_X(u)\to0$.
\end{theorem}

\begin{proof}
The decomposition \CGeqref{eq:dec0_k} holds termwise with weight $1/n$, $n=Lm+a$; the terms with $m=0$ contribute $O(1)$, and $\frac1{Lm+a}=\frac1{Lm}+O_L(m^{-2})$ for $m\ge1$. A block $(A_X,B_X]$ with $A>0$ has logarithmic mass $O(1)$; only blocks starting at $A=0$ need the following input. After a fixed shift making all intercepts positive, each block is $o(\log X)$ by \cite[Theorem~1.1, p.~2]{TaoTeravainen2018}, which for odd $k$ requires no non-degeneracy condition \cite[Remark~1.2, p.~2]{TaoTeravainen2018}. There are $O(1)$ blocks and $H_X\sim\log X$.
\end{proof}


\section{Three and more points: the centering step of the bias method}
\label{sec:three}

\subsection{The bias method for two points}
\label{subsec:bias}

We recall the structure of the proof of the binary corrected Elliott theorem \cite[Theorem~1.2, p.~4]{OpenAI2026}, Part~II of \cite{OpenAI2026}, specialised to $f_1=f_2=\lambda$. The parameters are \cite[(10.1)--(10.2), p.~44]{OpenAI2026}
\begin{equation}
\varepsilon=10^{-4},\quad \eta=\frac{\varepsilon}{100},\quad \kappa_0=\frac{400}{\eta},\quad A=e^{2\kappa_0},\quad c_*=\frac{\varepsilon}{100},\quad \theta=\frac14,
\label{eq:oai_params}
\end{equation}
\begin{equation}
\begin{aligned}
&\mathcal C=\{p:\ B^{1-\eta}<\log p\le B\},\qquad \mathcal Z=\{p:\ B^{1-\varepsilon}<\log p\le B^{1-\eta}\},\\
&\mathcal S=\mathcal C\cup\mathcal Z,\qquad P_0=e^{B^{1-\varepsilon}},
\end{aligned}
\label{eq:oai_primes}
\end{equation}
with $v_{\mathcal C}=\sum_{p\in\mathcal C}1/p=(\eta+o(1))\log B$ and $L_0=\prod_{p\in\mathcal C}(1+A/p)\prod_{p\in\mathcal Z}(1+1/p)$. An admissible family $\mathcal D$ consists of squarefree $\mathcal S$-products $d$ with $H<d\le\tau H$, $1\le H\le e^{C_0B}$, $\omega(d)\le J=\lceil C_0\log B\rceil$; $d_{\mathcal C}$ is the $\mathcal C$-part of $d$, and the cutoffs are $w_d(x)=\varphi((k_d(x)-v_{\mathcal C})/\sqrt{v_{\mathcal C}})$ with $k_d(x)=\#\{p\in\mathcal C:p\mid x,\ p\nmid d\}$ \cite[(10.3)--(10.4), (11.3), pp.~45--47]{OpenAI2026}. The proof has four steps:
\begin{enumerate}
\item a persistent bias $|S(N_j)|\ge\gamma N_j$ is assumed;
\item a divisor interval carries mass $\gg L_0/B$ \cite[Lemma~11.2, p.~47]{OpenAI2026}, and dilation $x=dm$ turns the bias into a lower bound $\gg_\gamma L_0/B$ for a raw sum with ordinary divisibility indicators \cite[Lemma~11.3, p.~48]{OpenAI2026};
\item centering the center-band indicators $\mathbf 1_{p\mid x}\mapsto\mathbf 1_{p\mid x}-\theta/p$ changes the raw sum by $\ll L_0(B^{-1-\varepsilon}+J/P_0)$ \cite[Proposition~11.4, p.~49]{OpenAI2026};
\item the centered sum is $\ll L_0B^{-1-c_*/2}$ by the divisor graph estimate \cite[Proposition~10.1, p.~45]{OpenAI2026}.
\end{enumerate}
Step~(3) is reduced \cite[\S12.1, pp.~50--51]{OpenAI2026} to the estimate
\begin{equation}
\limsup_{Y\to\infty}\Bigl|\frac1Y\sum_{z\le Y}b_1(z)\sum_{M<w\le\tau M}\frac{c_w}{w}\,b_2(z+hw)\Bigr|\ll B^{-1-\varepsilon},
\label{eq:oai124}
\end{equation}
uniformly in coefficients $|c_w|\le1$ supported on $P_0$-rough $w\in(M,\tau M]$, $P_0/\tau\le M\le e^{C_0B}$, and in the twists $b_i(n)=\lambda(n)e(t_i\#\{p\in\mathcal C:p\mid n,\ p\nmid u\}/\sqrt{v_{\mathcal C}})$. The proof of \CGeqref{eq:oai124} \cite[pp.~53--55]{OpenAI2026} uses of $b_1,b_2$ only their $1$-boundedness and the averaged short exponential sum bound \cite[Theorem~12.2 and (12.7), pp.~51--52]{OpenAI2026}
\begin{equation}
\limsup_{Y\to\infty}\ \sup_{\alpha\in\mathbb R/\mathbb Z}\ \frac1{DY}\int_0^Y\Bigl|\sum_{y<m\le y+D}b(m)e(\alpha m)\Bigr|\,dy\ll\frac{\log\log D}{\log D},
\label{eq:oai127}
\end{equation}
for one of the two functions, with the supremum outside the integral, together with the rough-number bounds of \cite[Lemma~12.3, p.~52]{OpenAI2026}: the sup-norm and fourth-moment bounds (12.9)--(12.10) for $\sum_wc_ww^{-1}e(\alpha hw)$, which rest on the count $\sum_{M<w\le\tau M,\ w\ P_0\text{-rough}}1/w\ll B^{-1+\varepsilon}\log B$ \cite[(12.13), p.~53]{OpenAI2026}.

For $k$ points $0=h_0<h_1<\dots<h_{k-1}$, the same expansion leads to
\begin{equation}
\mathcal E_k:=\limsup_{Y\to\infty}\Bigl|\frac1Y\sum_{z\le Y}b_0(z)\sum_{M<w\le\tau M}\frac{c_w}{w}\prod_{i=1}^{k-1}b_i(z+h_iw)\Bigr|\ll B^{-1-\varepsilon}.
\label{eq:E_k}
\end{equation}
For $f_i=\lambda$ the dilation coefficient of \cite[Lemma~11.3]{OpenAI2026} is $a_d=\lambda(d)^k$, and the expansion of \cite[\S12.1]{OpenAI2026}, with $d=uw$ and $x=uz$, produces the coefficients $c^{(u)}_w=a_{uw}(-\theta)^{\omega(w)}\mathbf 1[uw\in\mathcal D]$ together with the extracted factor $\prod_i\lambda(u)=\lambda(u)^k$. Since $w$ is squarefree, the effective coefficient is
\begin{equation}
\lambda(u)^k\,c^{(u)}_w=\lambda(w)^k(-\theta)^{\omega(w)}\mathbf 1[uw\in\mathcal D]=
\begin{cases}\theta^{\omega(w)}\mathbf 1[uw\in\mathcal D],&k\ \text{odd},\\ (-\theta)^{\omega(w)}\mathbf 1[uw\in\mathcal D],&k\ \text{even}.\end{cases}
\label{eq:cw_general}
\end{equation}
For odd $k$ every term of the expansion therefore carries a non-negative coefficient, independently of $u$.

\subsection{Barriers}
\label{subsec:barrier}

\begin{proposition}[Polynomial-phase barrier, all $k\ge3$]
\label{prop:barrier_poly}
Let $k\ge3$, $0=h_0<h_1<\dots<h_{k-1}$, $V:=\prod_{0\le j<l\le k-1}(h_l-h_j)$ and $R_i:=V\prod_{j\ne i}(h_i-h_j)^{-1}\in\mathbb Z\setminus\{0\}$. Let $\beta$ be irrational and $b_i(m):=e(\beta R_im^{k-1})$. Then:
\begin{enumerate}
\item for every integer $D\ge1$, uniformly in $\alpha\in\mathbb R$,
\begin{equation}
\limsup_{Y\to\infty}\frac1{DY}\int_0^Y\Bigl|\sum_{y<m\le y+D}b_i(m)e(\alpha m)\Bigr|\,dy\le D^{-1/2};
\label{eq:barrier_poly_127}
\end{equation}
\item $\prod_{i=0}^{k-1}b_i(z+h_iw)=e(\beta Vw^{k-1})$ for all $z,w$; hence with $c_w:=e(-\beta Vw^{k-1})\mathbf 1[w\ \text{is }P_0\text{-rough}]$,
\begin{equation}
\frac1Y\sum_{z\le Y}b_0(z)\sum_{M<w\le\tau M}\frac{c_w}{w}\prod_{i=1}^{k-1}b_i(z+h_iw)=\frac{\lfloor Y\rfloor}Y\sum_{\substack{M<w\le\tau M\\ w\ P_0\text{-rough}}}\frac1w .
\label{eq:barrier_poly_E}
\end{equation}
\end{enumerate}
For $P_0\le M\le P_0^2/\tau$, a range contained in that of \CGeqref{eq:oai124}, every $P_0$-rough integer of $(M,\tau M]$ is prime, and the right-hand side of \CGeqref{eq:barrier_poly_E} is $\sum_{M<p\le\tau M}1/p\ge(\tau-1)/(3\tau\log(\tau M))\gg_\tau B^{-1+\varepsilon}$ for large $B$, which exceeds $B^{-1-\varepsilon}$ by the factor $B^{2\varepsilon}$. Thus \CGeqref{eq:E_k} does not follow from \CGeqref{eq:oai127}, the rough support and the size of the coefficients.
\end{proposition}

\begin{proof}
By Lagrange interpolation, $\sum_ih_i^j\prod_{l\ne i}(h_i-h_l)^{-1}$ equals $0$ for $0\le j\le k-2$ and $1$ for $j=k-1$. Hence $\sum_iR_i(z+h_iw)^{k-1}=\sum_j\binom{k-1}{j}z^{k-1-j}w^j\sum_iR_ih_i^j=Vw^{k-1}$, which is (2). For (1), write $S_y(\alpha)$ for the inner sum; since $D$ is an integer it has exactly $D$ terms $m=n_0+t$, $1\le t\le D$, $n_0=\lfloor y\rfloor$. Expanding the square,
\begin{equation}
\begin{aligned}
\frac1Y\int_0^Y|S_y(\alpha)|^2dy=\sum_{1\le t,t'\le D}&e(\alpha(t-t'))\\
&\times\frac1Y\int_0^Ye\bigl(\beta R_i\bigl((n_0+t)^{k-1}-(n_0+t')^{k-1}\bigr)\bigr)dy .
\end{aligned}
\label{eq:barrier_cs}
\end{equation}
For $t\ne t'$, $n\mapsto\beta R_i((n+t)^{k-1}-(n+t')^{k-1})$ is a polynomial of degree $k-2\ge1$ with irrational leading coefficient $\beta R_i(k-1)(t-t')$, so by Weyl's equidistribution theorem its average over $n_0\le Y$ tends to $0$. There are finitely many pairs, and $\alpha$ enters only through unimodular constants, so the left-hand side of \CGeqref{eq:barrier_cs} tends to $D$ uniformly in $\alpha$. Cauchy--Schwarz gives $\limsup_Y\frac1Y\int_0^Y|S_y|\,dy\le D^{1/2}$. The lower bound for the prime sum is the prime number theorem, with $\log(\tau M)\le2B^{1-\varepsilon}$.
\end{proof}

For $f_i=\lambda$ the coefficients \CGeqref{eq:cw_general} are not arbitrary: they are $(\pm\theta)^{\omega(w)}$ on the admissible $w$, and for odd $k$ they are non-negative; Proposition~\CGref{prop:barrier_poly} does not model them. For odd $k$ the polynomial phases of Proposition~\CGref{prop:barrier_poly} then no longer produce a constant product. For $k=3$ the barrier survives with phases of vanishing curvature.

\begin{proposition}[Barrier for non-negative coefficients, $k=3$]
\label{prop:barrier_pos}
Let $0<h_1<h_2$, $\rho:=h_1h_2(h_1-h_2)$, $\tau\in(1,2)$, $H,D^*\ge1$, and let $q$ be a prime with
\begin{equation}
q\ge\max\bigl(2h_2+1,\ 7h_2(D^*+1)^3,\ 6|\rho|\tau^2H^2\bigr).
\label{eq:q_choice}
\end{equation}
Put $b_0(m)=e((h_1-h_2)m^2/q)$, $b_1(m)=e(h_2m^2/q)$, $b_2(m)=e(-h_1m^2/q)$. Then:
\begin{enumerate}
\item for integers $1\le D\le D^*$ and every $\alpha$,
\begin{equation}
\limsup_{Y\to\infty}\frac1{DY}\int_0^Y\Bigl|\sum_{y<m\le y+D}b_i(m)e(\alpha m)\Bigr|dy\le\frac{3+\log q}{D};
\label{eq:barrier_pos_127}
\end{equation}
\item for $M\le H$ and every $c_w\ge0$ supported on $(M,\tau M]$,
\begin{equation}
\begin{aligned}
\liminf_{Y\to\infty}\ \mathrm{Re}\,\frac1Y\sum_{z\le Y}b_0(z)&\sum_{M<w\le\tau M}\frac{c_w}{w}b_1(z+h_1w)b_2(z+h_2w)\\
&\ge\ \frac12\sum_{M<w\le\tau M}\frac{c_w}{w}.
\end{aligned}
\label{eq:barrier_pos}
\end{equation}
\end{enumerate}
\end{proposition}

\begin{proof}
(2) $(h_1-h_2)z^2+h_2(z+h_1w)^2-h_1(z+h_2w)^2=\rho w^2$, so the product equals $e(\rho w^2/q)$, and $|\rho|w^2/q\le|\rho|\tau^2H^2/q\le\frac16$ gives $\mathrm{Re}\,e(\rho w^2/q)\ge\frac12$. (1) Let $r\in\{h_1-h_2,h_2,-h_1\}$, $n_0=\lfloor y\rfloor$, and $L\le D+1$ the number of terms. Then $r(n_0+t)^2/q+\alpha(n_0+t)=\mathrm{const}+\vartheta t+rt^2/q$ with $\vartheta=\alpha+2rn_0/q$, and $\sum_{t\le L}|e(rt^2/q)-1|\le2\pi h_2L^3/q\le1$. So $|S_y|\le\min(L,\|\vartheta\|^{-1}/2)+1$. As $n_0$ runs over a period modulo $q$, $2rn_0$ runs over all residues ($q\nmid2r$). The points $\alpha+j/q$ have spacing $1/q$, at most two of them lie at distance $<1/q$ from $\mathbb Z$, and at most two lie at distance in $[j'/q,(j'+1)/q)$ for each $j'\ge1$; hence the average of $\min(L,\|\alpha+j/q\|^{-1}/2)$ over $j$ is at most $2L/q+1+\log q\le2+\log q$.
\end{proof}

The functions of Proposition~\CGref{prop:barrier_pos} are $q$-periodic, so the bound in (1) holds only for windows $D\le D^*$; this is all that matters. In \cite[pp.~54--55]{OpenAI2026} the bound \CGeqref{eq:oai127} is used only at the two window lengths $D_1\asymp M$ and $D_2\asymp_hD_1$ fixed before $Y\to\infty$. Choosing $D^*\ge(1+2h_2)\lceil H\rceil$ gives $q\asymp_{h}H^3$, $\log q\ll B$, and for $P_0/\tau\le D\le D^*$ the bound in (1) is $O(Be^{-B^{1-\varepsilon}})$, far below $\log\log D/\log D$. So the functions of Proposition~\CGref{prop:barrier_pos} satisfy \CGeqref{eq:oai127} at every window length used in \cite[\S12]{OpenAI2026}. Their product along the pattern is $\ge\frac12$ in real part, uniformly in $u$ when $q$ is chosen for $H$, so for odd $k=3$ no cancellation between different $u$ occurs either. The functions are locally linear phases with slowly varying frequency: on each window the frequency is $\alpha+2rn_0/q$. Any extension of Step~(3) to three points must therefore use information on $b_i$ that excludes such local linear phases, which \CGeqref{eq:oai127} does not.

\subsection{Local Gowers norms control the centering error}
\label{subsec:local_gowers}

For a prime $N$ and $g:\mathbb Z_N\to\mathbb C$ let $\mathbb E$ denote averages over $\mathbb Z_N$, $\Delta_ag(x):=g(x)\overline{g(x+a)}$, $\|g\|_{U^1}:=|\mathbb Eg|$, and
\begin{equation}
\|g\|_{U^{s}}^{2^{s}}:=\mathbb E_{a}\,\|\Delta_ag\|_{U^{s-1}}^{2^{s-1}}\qquad(s\ge2),
\label{eq:gowers}
\end{equation}
the Gowers norms \cite{Gowers2001}. In particular $\|g\|_{U^2}^4=\sum_\xi|\widehat g(\xi)|^4\ge\max_\xi|\widehat g(\xi)|^4$, with $\widehat g(\xi)=\mathbb E_xg(x)e(-x\xi/N)$.

\begin{definition}[Local Gowers norm]
\label{def:local_gowers}
For an integer $D\ge2$ let $N_D$ be the least prime $>4D$. For $b:\mathbb N\to\mathbb C$ and $y\ge0$ define $G_y:\mathbb Z_{N_D}\to\mathbb C$ by $G_y(x)=b(\lfloor y\rfloor+x)$ for $x\in\{1,\dots,2D+1\}$ and $G_y(x)=0$ for the other residues, and put $\|b\|_{U^s(y;D)}:=\|G_y\|_{U^s(\mathbb Z_{N_D})}$.
\end{definition}

\begin{lemma}[Generalised von Neumann inequality]
\label{lem:vonneumann}
Let $k\ge2$, $0=h_0<h_1<\dots<h_{k-1}$, $N>h_{k-1}$ prime, and $g_0,\dots,g_{k-1}:\mathbb Z_N\to\mathbb C$ with $|g_i|\le1$. Then
\begin{equation}
\Bigl|\mathbb E_{z,w\in\mathbb Z_N}\prod_{i=0}^{k-1}g_i(z+h_iw)\Bigr|\le\|g_{k-1}\|_{U^{k-1}}.
\label{eq:vonneumann}
\end{equation}
\end{lemma}

\begin{proof}
For $k=2$ the average factorises as $\mathbb E g_0\cdot\mathbb Eg_1$, of modulus $\le\|g_1\|_{U^1}$. For $k\ge3$, Cauchy--Schwarz in $z$ and $|g_0|\le1$ give
\begin{equation}
\begin{aligned}
\Bigl|\mathbb E_{z,w}\prod_{i=0}^{k-1}g_i(z+h_iw)\Bigr|^2&\le\mathbb E_{z}\Bigl|\mathbb E_w\prod_{i=1}^{k-1}g_i(z+h_iw)\Bigr|^2\\
&=\mathbb E_t\,\mathbb E_{z,w}\prod_{i=1}^{k-1}(\Delta_{h_it}g_i)(z+(h_i-h_1)w),
\end{aligned}
\label{eq:vn_step}
\end{equation}
after writing $w'=w+t$ and replacing $z$ by $z-h_1w$. The shifts $0<h_2-h_1<\dots<h_{k-1}-h_1$ are distinct and $<N$, so by induction the inner average is at most $\|\Delta_{h_{k-1}t}g_{k-1}\|_{U^{k-2}}$ in modulus. Since $t\mapsto h_{k-1}t$ is a bijection of $\mathbb Z_N$, Hölder and \CGeqref{eq:gowers} give $\mathbb E_t\|\Delta_{h_{k-1}t}g_{k-1}\|_{U^{k-2}}\le(\mathbb E_s\|\Delta_sg_{k-1}\|_{U^{k-2}}^{2^{k-2}})^{2^{2-k}}=\|g_{k-1}\|_{U^{k-1}}^2$.
\end{proof}

\begin{proposition}[Local $U^k$ control of the $k$-point centering error]
\label{prop:local_control}
Let $k\ge3$, $0=h_0<h_1<\dots<h_{k-1}$, $\tau\in(1,2)$, $M\ge1$, $\mathcal J=(M,\tau M]$, $|c_w|\le1$ for $w\in\mathcal J$, $\varpi:=\sum_{w\in\mathcal J}|c_w|/w$, and $D:=\lceil h_{k-1}\tau M\rceil$. There is $C=C(h_{k-1})$ such that for all $b_0,\dots,b_{k-1}:\mathbb N\to\mathbb C$ with $|b_i|\le1$ and all $Y\ge1$,
\begin{equation}
\begin{aligned}
&\frac1Y\Bigl|\sum_{z\le Y}b_0(z)\sum_{w\in\mathcal J}\frac{c_w}w\prod_{i=1}^{k-1}b_i(z+h_iw)\Bigr|\\
&\qquad\le C\,\bigl(\varpi\log(3M)\bigr)^{1/2}\,\frac1Y\int_0^Y\|b_{k-1}\|_{U^k(y;D)}\,dy\\
&\qquad\quad+\frac{(2D+1)\varpi}{Y}.
\end{aligned}
\label{eq:local_control}
\end{equation}
No hypothesis on the support of $c_w$ or on $b_0,\dots,b_{k-2}$ is needed.
\end{proposition}

\begin{proof}
Write $g(z)$ for the summand in $z$, so $|g|\le\varpi$, and for $y\ge0$ let $\mathcal W_y=(y,y+D]\cap\mathbb Z$ and $S_y=\sum_{z\in\mathcal W_y}g(z)$. Each integer $z\in[D,Y]$ lies in $\mathcal W_y$ for a set of $y\in[0,Y]$ of measure $D$, and every other $z$ for a set of measure in $[0,D]$; hence
\begin{equation}
\Bigl|\sum_{z\le Y}g(z)\Bigr|\le\frac1D\int_0^Y|S_y|\,dy+(2D+1)\varpi .
\label{eq:windows}
\end{equation}
\emph{Cauchy--Schwarz in $w$.} With $X_{y,w}:=\sum_{z\in\mathcal W_y}b_0(z)\prod_{i\ge1}b_i(z+h_iw)$ and $|c_w|/w\le1/M$,
\begin{equation}
|S_y|^2\le\varpi\sum_{w\in\mathcal J}\frac{|c_w|}{w}|X_{y,w}|^2\le\frac{\varpi}{M}\sum_{w\in\mathcal J}|X_{y,w}|^2 .
\label{eq:cs_w}
\end{equation}
\emph{Differencing.} Put $n_0=\lfloor y\rfloor$, so that $\mathcal W_y-n_0=\{1,\dots,D\}$, and $G_i(x):=b_i(n_0+x)\mathbf 1_{1\le x\le 2D+1}$ and $F_0(x):=b_0(n_0+x)\mathbf 1_{n_0+x\in\mathcal W_y}$. Writing $z'=z+s$ in $|X_{y,w}|^2$,
\begin{equation}
\sum_{w\in\mathcal J}|X_{y,w}|^2=\sum_{|s|<D}T_s,\qquad T_s:=\sum_{x}\sum_{w\in\mathcal J}(\Delta_sF_0)(x)\prod_{i=1}^{k-1}(\Delta_sG_i)(x+h_iw),
\label{eq:Tk}
\end{equation}
since for $x,x+s\in\mathcal W_y-n_0\subseteq\{1,\dots,D+1\}$ and $w\in\mathcal J$ all arguments $x+h_iw$, $x+s+h_iw$ lie in $\{1,\dots,2D+1\}$ (here $h_{k-1}w\le h_{k-1}\tau M\le D$).

\emph{Transfer to $\mathbb Z_N$.} Let $N=N_D\in(4D,8D]$ (Bertrand). Read $F_0,G_i$ as functions on $\mathbb Z_N$ supported on the residues of $\{1,\dots,D+1\}$ and $\{1,\dots,2D+1\}$, and take $\Delta_s$ in $\mathbb Z_N$. For $|s|<D$, the arguments $x+s$ and $x+h_iw+s$ with $x\le D+1$ lie in $\{2-D,\dots,3D\}$; those outside $\{1,\dots,3D\}$ lie in $\{2-D,\dots,0\}$, whose residues are $\ge N-D+2>2D+1$; so $\Delta_sF_0$ is unchanged, and for $w\in\mathcal J$ and $x$ in the support of $\Delta_sF_0$ the values $(\Delta_sG_i)(x+h_iw)$ computed in $\mathbb Z_N$ equal the integer ones. Thus $T_s=\sum_{x,w\in\mathbb Z_N}\mathbf 1_{\mathcal J}(w)(\Delta_sF_0)(x)\prod_{i\ge1}(\Delta_sG_i)(x+h_iw)$, where $\mathcal J$ is read in $\mathbb Z_N$ ($\tau M<N$). Expand $\mathbf 1_{\mathcal J}(w)=\sum_\xi\widehat{\mathbf 1_{\mathcal J}}(\xi)e(\xi w/N)$ with $\sum_\xi|\widehat{\mathbf 1_{\mathcal J}}(\xi)|\le2+\log N$, and write $e(\xi w/N)=e(\xi\bar h_1(x+h_1w)/N)\,e(-\xi\bar h_1x/N)$ with $\bar h_1h_1\equiv1$. The two linear phases are absorbed into $\Delta_sF_0$ and $\Delta_sG_1$ without changing their sup norms; since $k\ge3$, $G_1\ne G_{k-1}$ and the function measured in Lemma~\CGref{lem:vonneumann} is untouched, so
\begin{equation}
|T_s|\le N^2(2+\log N)\,\|\Delta_sG_{k-1}\|_{U^{k-1}(\mathbb Z_N)} .
\label{eq:Tk_bound}
\end{equation}
\emph{Summing over $s$.} By Hölder and \CGeqref{eq:gowers},
\begin{equation}
\sum_{|s|<D}\|\Delta_sG_{k-1}\|_{U^{k-1}}\le N^{1-2^{1-k}}\Bigl(\sum_{s\in\mathbb Z_N}\|\Delta_sG_{k-1}\|^{2^{k-1}}_{U^{k-1}}\Bigr)^{2^{1-k}}=N\,\|G_{k-1}\|^2_{U^k}.
\label{eq:sum_k}
\end{equation}
Combining \CGeqref{eq:cs_w}--\CGeqref{eq:sum_k}, $|S_y|\le(\varpi N^3(2+\log N)/M)^{1/2}\|G_{k-1}\|_{U^k}$. Since $\|G_{k-1}\|_{U^k}=\|b_{k-1}\|_{U^k(y;D)}$, $N\le8D$ and $D\le3h_{k-1}M$, one has $N^3/(MD^2)\le1536h_{k-1}$ and $\log N\ll_{h_{k-1}}\log(3M)$. Insert in \CGeqref{eq:windows} and divide by $Y$.
\end{proof}

\subsection{A conditional extension of the bias method}
\label{subsec:conditional}

\begin{hypothesis}[$\mathrm{LU}_k(\varrho,K)$: quantitative local Gowers uniformity of $\lambda$]
\label{hyp:LU}
Let $k\ge3$. For every integer $D\ge2$, every finite set of primes $\mathcal P$ and every multiplicative $b:\mathbb N\to\mathbb C$ with $|b|\le1$ and $b(p^j)=\lambda(p^j)$ for all $p\notin\mathcal P$, $j\ge1$,
\begin{equation}
\limsup_{Y\to\infty}\frac1Y\int_0^Y\|b\|_{U^k(y;D)}\,dy\le K\Bigl(2+\sum_{p\in\mathcal P}\frac1p\Bigr)^K(\log D)^{-\varrho}.
\label{eq:LU}
\end{equation}
\end{hypothesis}

\begin{theorem}[$k$-point centering under $\mathrm{LU}_k$]
\label{thm:centering_k}
Let $k\ge3$ and assume $\mathrm{LU}_k(\varrho,K)$ for some $K$ and some $\varrho>1+3\varepsilon$, with $\varepsilon=10^{-4}$ as in \CGeqref{eq:oai_params}. Then \CGeqref{eq:E_k} holds for the twists $b_i$ of \cite[(12.3), p.~51]{OpenAI2026}, any $1$-bounded $b_0,\dots,b_{k-2}$, and all coefficients $|c_w|\le1$ supported on $P_0$-rough integers of $(M,\tau M]$ with $P_0/\tau\le M\le e^{C_0B}$, uniformly, for all sufficiently large $B$.
\end{theorem}

\begin{proof}
The twist $b_{k-1}$ agrees with $\lambda$ at all prime powers of primes outside $\mathcal C$, and $\sum_{p\in\mathcal C}1/p=v_{\mathcal C}\le\log B$. By \cite[(12.13), p.~53]{OpenAI2026}, $\varpi\ll B^{-1+\varepsilon}\log B$; also $\log(3M)\ll B$ and $\log D\ge\log M\ge B^{1-\varepsilon}-1$. Proposition~\CGref{prop:local_control} and \CGeqref{eq:LU} give
\begin{equation}
\mathcal E_k\ll\bigl(B^{-1+\varepsilon}(\log B)\,B\bigr)^{1/2}(2+\log B)^K B^{-\varrho(1-\varepsilon)}\ll B^{\varepsilon/2}(\log B)^{K+1/2}B^{-\varrho(1-\varepsilon)},
\label{eq:Ek_bound}
\end{equation}
and $\varrho(1-\varepsilon)>1+\frac32\varepsilon$ when $\varrho>1+3\varepsilon$. The boundary term $(2D+1)\varpi/Y$ of \CGeqref{eq:local_control} tends to $0$ because $D$ is fixed before $Y\to\infty$.
\end{proof}

For the graph step we formulate the $k$-linear analogue of \cite[Proposition~10.1, p.~45]{OpenAI2026}. With the notation of Subsection~\CGref{subsec:bias} put
\begin{equation}
g^{(k)}_d(x):=A^{\omega_{\mathcal C}(d)}\,\mathbf 1_{d_{\mathcal C}\mid x}\prod_{\substack{p\mid d\\ p\in\mathcal Z}}\Bigl(\mathbf 1_{p\mid x}-\frac\theta p\Bigr)\prod_{i=0}^{k-1}w_d(x+h_id).
\label{eq:gk}
\end{equation}

\begin{hypothesis}[$\mathrm{G}_k(h)$: $k$-linear divisor graph estimate]
\label{hyp:G}
Fix $0=h_0<h_1<\dots<h_{k-1}$, $1<\tau<2$, $C_0\ge1$, $T>0$. For every sufficiently large $B$, every admissible family $\mathcal D$, all cutoffs $w_d$ as in \cite[(10.4)]{OpenAI2026}, all $|a_d|\le1$ and all $F_0,\dots,F_{k-1}:\mathbb N\to\mathbb D$,
\begin{equation}
\limsup_{X\to\infty}\frac1X\Bigl|\sum_{x\le X}\sum_{d\in\mathcal D}a_d\prod_{i=0}^{k-1}F_i(x+h_id)\,g^{(k)}_d(x)\Bigr|\ll L_0B^{-1-c_*/2},
\label{eq:Gk}
\end{equation}
with implied constant depending only on $h,\tau,C_0,T$ and \CGeqref{eq:oai_params}.
\end{hypothesis}

For $k=2$, \CGeqref{eq:Gk} is \cite[Proposition~10.1]{OpenAI2026}, proved there through high traces of a self-adjoint graph operator \cite[\S\S13--17, pp.~55--81]{OpenAI2026}; for $k\ge3$ the natural object is a $k$-uniform hypergraph with hyperedges $\{x+h_id\}_i$, and the trace method has no direct analogue.

\begin{theorem}[Conditional $k$-point Chowla in progressions]
\label{thm:conditional_k}
Let $k\ge3$ and $0=h_0<\dots<h_{k-1}$. Assume $\mathrm{LU}_k(\varrho,K)$ with $\varrho>1+3\varepsilon$ and $\mathrm{G}_k(h)$. Then for all fixed $l\ge1$ and $b\in\mathbb Z$,
\begin{equation}
\frac1N\sum_{\substack{m\le N\\ m\equiv b\ (\mathrm{mod}\ l)}}\prod_{i=0}^{k-1}\lambda(m+h_i)\longrightarrow0\qquad(N\to\infty).
\label{eq:conditional_k}
\end{equation}
\end{theorem}

\begin{proof}
Suppose the contrary: there are $\gamma\in(0,1]$ and $N_j\to\infty$ with $|S(N_j)|\ge\gamma N_j$, where $S(N)$ is the sum in \CGeqref{eq:conditional_k}. We follow \cite[\S11]{OpenAI2026} with three changes.

\emph{(a) Divisors $\equiv1\pmod l$.} For large $B$ every prime of $\mathcal S$ exceeds $l$. Under the divisor law of \cite[proof of Lemma~11.2]{OpenAI2026} each $p\in\mathcal S$ is included independently with probability $A_a/(p+A_a)$; for a non-principal character $\chi$ modulo $l$,
\begin{equation}
\begin{aligned}
|\mathbb E_{\mathrm{div}}\chi(d)|&=\prod_{p\in\mathcal S}\Bigl|\frac{1+A_a\chi(p)/p}{1+A_a/p}\Bigr|\\
&\le\exp\Bigl(-\sum_{p\in\mathcal Z}\frac{1-\mathrm{Re}\,\chi(p)}{p}+O\Bigl(\sum_{p>P_0}p^{-2}\Bigr)\Bigr)\to0,
\end{aligned}
\label{eq:char_div}
\end{equation}
because $\sum_{p\in\mathcal Z}1/p=v_{\mathcal Z}\to\infty$ while $\sum_{p\in\mathcal Z}\chi(p)/p\to0$ by the prime number theorem in progressions modulo $l$. Hence $\mathbb P_{\mathrm{div}}(d\equiv1\ (\mathrm{mod}\ l))\to1/\varphi(l)$. Choosing $C_0$ so that $\mathbb P_{\mathrm{div}}(\omega(d)>J)$ and $\mathbb P_{\mathrm{div}}(\log d>C_0B)$ are each $\le1/(8\varphi(l))$, the argument of \cite[Lemma~11.2]{OpenAI2026} gives an interval $(H,\tau H]$ and a family $\mathcal D$ of divisors $d\equiv1\pmod l$ with $\sum_{d\in\mathcal D}A^{\omega_{\mathcal C}(d)}/d\ge c\,L_0/B$, $c=c(l,\tau,C_0)>0$; the condition $|f_1f_2(d)|\ge\frac12$ is void for $\lambda$.

\emph{(b) Raw lower bound with $k$ cutoffs.} Put $a_d=\lambda(d)^k$, $X_j=HN_j$, and
\begin{equation}
R_B(X):=\sum_{\substack{x\le X\\ x\equiv b\ (l)}}\sum_{d\in\mathcal D}a_dA^{\omega_{\mathcal C}(d)}\prod_{i=0}^{k-1}\lambda(x+h_id)w_d(x+h_id)\,\mathbf 1_{d\mid x}.
\label{eq:RB_k}
\end{equation}
For $x=dm$, $d\equiv1\pmod l$ gives $x\equiv b\iff m\equiv b\pmod l$, and complete multiplicativity gives $a_d\prod_i\lambda(d(m+h_i))=\prod_i\lambda(m+h_i)$ with no coprimality condition. As in \cite[proof of Lemma~11.3]{OpenAI2026}, using that the primes of $\mathcal C$ exceed $l$ (so that, by the Chinese remainder theorem, divisibility by them is independent of the class $m\bmod l$), Chebyshev's inequality bounds the proportion of $m$ on which some of the $k$ cutoffs differs from $1$ by $16k/T^2$, and the comparison of $X_j/d\in[N_j/\tau,N_j]$ with $N_j$ costs $\tau-1$. Choosing $\tau-1\le\gamma/64$ and $16k/T^2\le\gamma/64$ gives $\liminf_j|R_B(X_j)|/X_j\ge c_\gamma L_0/B$.

\emph{(c) Centering.} Let $C_B(X)$ be \CGeqref{eq:RB_k} with the factor $A^{\omega_{\mathcal C}(d)}\prod_iw_d(x+h_id)\mathbf 1_{d\mid x}$ replaced by $g^{(k)}_d(x)$. The expansion of \cite[\S12.1]{OpenAI2026} applies verbatim: $d=uw$, $x=uz$, and $\prod_i\lambda(u(z+h_iw))$ equals $\lambda(u)^k\prod_i\lambda(z+h_iw)$; the residue condition $uz\equiv b$ is absorbed into $b_0$, and the $k$ Fourier inversions of $\varphi$ cost $\|\widehat\varphi\|_1^k$. By Theorem~\CGref{thm:centering_k} and $\sum_uA^{\omega_{\mathcal C}(u)}/u=L_0$, $\limsup_X|R_B(X)-C_B(X)|/X\ll L_0B^{-1-\varepsilon}$.

\emph{(d) Contradiction.} Apply $\mathrm{G}_k(h)$ with $F_0(x)=\mathbf 1[x\equiv b\ (\mathrm{mod}\ l)]\lambda(x)$ and $F_i=\lambda$ for $i\ge1$. It gives $\limsup_X|C_B(X)|/X\ll L_0B^{-1-c_*/2}$. Hence $c_\gamma L_0/B\le C L_0(B^{-1-c_*/2}+B^{-1-\varepsilon})$, which fails for large $B$. The order of limits is that of \cite[p.~49]{OpenAI2026}: $\gamma,\tau,T,\varphi,C_0$ first, then $B$ fixed while $j\to\infty$, then $B\to\infty$.
\end{proof}

\begin{corollary}[The $k$-point family under $\mathrm{LU}$ and $\mathrm{G}$]
\label{cor:family_conditional}
Let $K_0\ge3$. Assume, for every $3\le k\le K_0$, $\mathrm{LU}_k(\varrho,K)$ with $\varrho>1+3\varepsilon$ and $\mathrm{G}_k(h)$ for every $h$. Then $\mathrm{Ch}_{k'}$ holds for $3\le k'\le K_0$, and \CGeqref{eq:limit_k} holds for every $k\le K_0$, every rational resonance and every $c$. In particular \CGeqref{eq:chowla013} holds.
\end{corollary}

\begin{proof}
Let $\alpha_jm+\beta_j$, $0\le j<k'$, be pairwise non-proportional, $Q=\prod_j\alpha_j$ and $N=Qm$. Then $(Q/\alpha_j)(\alpha_jm+\beta_j)=N+s_j$ with $s_j=(Q/\alpha_j)\beta_j$, pairwise distinct by non-proportionality, and $\prod_j\lambda(\alpha_jm+\beta_j)=\pm\prod_j\lambda(N+s_j)$. Ordering the $s_j$ and putting $M=N+\min_js_j$ turns the sum into a sum of $\prod_i\lambda(M+h_i)$ over $M\equiv\min_js_j\pmod Q$ with $0=h_0<\dots<h_{k'-1}$, which is $o(Y)$ by Theorem~\CGref{thm:conditional_k}. Theorem~\CGref{thm:suf_k} gives \CGeqref{eq:limit_k}.
\end{proof}

\begin{remark}[Channel form at a prime]
\label{rem:channels}
Fix a prime $p$ and shifts $H=\{h_0,\dots,h_{k-1}\}$, and write $C_H(X)=\frac1X\sum_{n\le X}\prod_r\lambda(n+h_r)$. Complete multiplicativity gives $\lambda(n+h_r)=(-1)^{v_p(n+h_r)}\lambda(U_{p,r}(n))$, where $U_{p,r}(n)$ is the $p$-free part of $n+h_r$, hence $\prod_r\lambda(n+h_r)=\eta_{p,H}(n)b_{p,H}(n)$ with $\eta_{p,H}(n)=(-1)^{\sum_rv_p(n+h_r)}$ and $b_{p,H}(n)=\prod_r\lambda(U_{p,r}(n))$. The finite-conductor encoding of \cite{CasanaPadic} (compendium entry 48.0.35) states that for every $J\ge1$
\begin{equation}
C_H(X)=\mu_p(H)T_{p,H,0}(X)+\sum_{j=1}^J\sum_{\substack{u\bmod p^j\\ p\nmid u}}\widehat{\Delta}_{p,H,j}(u)\,T_{p,H,j,u}(X)+O\Bigl(kp^{-J}+\frac kX\Bigr),
\label{eq:channels}
\end{equation}
with channels $T_{p,H,j,u}(X)=\frac1X\sum_{n\le X}b_{p,H}(n)e(un/p^j)$, and that $C_H(X)\to0$ as soon as every channel tends to $0$ (entry 48.0.36). Since $b_{p,H}=\eta_{p,H}\prod_r\lambda(\cdot+h_r)$ and $\eta_{p,H}$ coincides with a function of period $p^{J'}$ outside a set of density $O(kp^{-J'})$, each channel is, up to $O(kp^{-J'})$, a finite combination of the progression sums \CGeqref{eq:conditional_k} modulo $p^{\max(j,J')}$. Under the hypotheses of Theorem~\CGref{thm:conditional_k} every channel therefore tends to $0$. The encoding \CGeqref{eq:channels} is an exact reformulation and adds no cancellation of its own; the input it needs is exactly Theorem~\CGref{thm:conditional_k} in progressions to prime-power moduli.
\end{remark}

\begin{remark}[Status of the two hypotheses]
\label{rem:status}
Proposition~\CGref{prop:local_control}, Lemma~\CGref{lem:vonneumann} and Propositions~\CGref{prop:barrier_poly}--\CGref{prop:barrier_pos} are unconditional; Theorem~\CGref{thm:centering_k} is the implication $\mathrm{LU}_k\Rightarrow$\CGeqref{eq:E_k}, and Theorem~\CGref{thm:conditional_k} is conditional on both hypotheses. At fixed $D$, $\mathrm{LU}_k$ is a statement about averaged $2^k$-point correlations of $b$ at shifts $\le3D$ (expand $\|G_y\|_{U^k}^{2^k}$); it is therefore not a statement obviously weaker than $k$-point Chowla, and Theorem~\CGref{thm:conditional_k} should be read as isolating exactly which uniformity the bias method needs. $\mathrm{LU}_k$ concerns windows of fixed length $D$ while $Y\to\infty$, with an ordinary average in $y$ and the rate $(\log D)^{-\varrho}$. Matomäki--Radziwi{\l}{\l}--Tao proved local Fourier uniformity, with the supremum over frequencies inside the integral, for windows $H\ge X^\theta$ with any fixed $\theta>0$ \cite[Theorem~1.2, p.~4]{MRT2020}; Tao showed that the logarithmically averaged Chowla conjecture is equivalent to a logarithmically averaged local Gowers uniformity of $\lambda$ \cite{Tao2017}. Windows of fixed length, ordinary averages and a rate, as in \CGeqref{eq:LU}, are not covered by either. Proposition~\CGref{prop:barrier_pos} shows that, for $k=3$, some information excluding local linear phases is necessary for Step~(3); Theorem~\CGref{thm:centering_k} shows that the quantitative, uniform local $U^3$ hypothesis $\mathrm{LU}_3(\varrho,K)$ is sufficient. Whether local $U^2$ uniformity suffices for the specific non-negative coefficients of \CGeqref{eq:cw_general} is open. $\mathrm{G}_k$ for $k\ge3$ is open.
\end{remark}

\section{Numerical verification}
\label{sec:numerics}

All values below were regenerated from the supplied signed-int8 Liouville
prefix using the supplied Python scripts. They are finite checks and are not
used as asymptotic rate assumptions.

\begin{experiment}[Fixed rational resonances]
At $X=2\cdot10^6$:
\begin{equation}
\begin{array}{c|r|r|r}
p/q & C_X/X & \lambda(pq)/q & (C_X-X\lambda(pq)/q)/\sqrt X\\ \hline
3/2 &  0.4999230 &  0.5000000 & -0.109\\
5/2 &  0.4995180 &  0.5000000 & -0.682\\
7/2 &  0.4996260 &  0.5000000 & -0.529\\
2/3 &  0.3326445 &  0.3333333 & -0.974\\
5/3 &  0.3324220 &  0.3333333 & -1.289\\
4/3 & -0.3322400 & -0.3333333 &  1.546\\
5/4 & -0.2492080 & -0.2500000 &  1.120
\end{array}
\label{eq:o34}
\end{equation}
The $2/3$ computation omits the unique initial term for which
$\lfloor(2/3)n\rfloor=0$, as required by the definition.
\end{experiment}

\begin{figure}[htbp]
\centering
\includegraphics[width=.88\textwidth]{rational_resonances_X2000000.png}
\caption{Measured rational resonances at $X=2\cdot10^6$ and the conditional
natural main term.}
\label{fig:rational}
\end{figure}

\begin{experiment}[Natural microscopic profile]
For $p/q=3/2$ and $X=2\cdot10^6$:
\begin{equation}
\begin{array}{c|r|r}
c & C_X/X & \Psi_{3,2}(c)\\ \hline
-0.5 & -0.000265 & 0\\
-0.1 & -0.000321 & 0\\
0.0 & 0.499923 & 0.500000\\
0.3 & 0.499923 & 0.500000\\
0.6 & 0.499904 & 0.500000\\
0.8 & 0.415884 & 0.416667\\
1.2 & 0.276538 & 0.277778\\
2.0 & 0.166553 & 0.166667
\end{array}
\label{eq:o35}
\end{equation}
\end{experiment}

\begin{figure}[htbp]
\centering
\includegraphics[width=.88\textwidth]{rational_local_profile_3_over_2.png}
\caption{Natural microscopic profile around $\log(3/2)$.}
\label{fig:profile}
\end{figure}

\begin{experiment}[Integer endpoint, left interval, and majority]
For $m=2$, $X=10^6$, $c=1/2$,
\begin{equation}
Xm(e^{c/X}-1)=1.000000250000041666\ldots,\qquad C_X=-999998.
\label{eq:o36}
\end{equation}
For $m=2$, $X=10^7$, both $c=-0.01$ and $c=-0.40$ give exactly
$C_X=-778$. Two finite checks of Theorem~\CGref{thm:naturalmajority} are
\begin{equation}
\begin{array}{c|r|r|r}
c & \lambda(2)C_X/X & (2|I_{0,X}|-X)/X &
2\min(1,(2c)^{-1})-1\\ \hline
0.60 & 0.833072800 & 0.666666600 & 0.666666667\\
0.90 & 0.555258200 & 0.111111000 & 0.111111111
\end{array}
\label{eq:o37}
\end{equation}
The middle column is the exact finite-$X$ signed lower bound used in the
proof; the last column is its asymptotic limit.
\end{experiment}

\begin{experiment}[Logarithmic mesoscopic majority]
For $m=2$, $X=10^7$ and $c_X=X^\theta$, let $B_X$ denote the exact
$H_X$-normalized harmonic mass of the zero-correction block. The finite bound
in the proof is $2B_X-1$, not its limiting value $1-2\theta$. The regenerated
values are
\begin{equation}
\begin{array}{c|r|r|r|r}
\theta & c_X & \lambda(2)C_X^{\log} & 2B_X-1 & 1-2\theta\\ \hline
0.00 & 1.0000    & 0.958453098 & 0.916965031 & 1.000000000\\
0.20 & 25.1189   & 0.765504459 & 0.530794400 & 0.600000000\\
0.40 & 630.9573  & 0.573145613 & 0.144624375 & 0.200000000\\
0.45 & 1412.5375 & 0.523535345 & 0.048073466 & 0.100000000
\end{array}
\label{eq:o38}
\end{equation}
Thus the finite measured values satisfy the exact finite-$X$ majority bound;
the final column is displayed only as the asymptotic lower bound.
\end{experiment}

\begin{figure}[htbp]
\centering
\includegraphics[width=.88\textwidth]{logarithmic_mesoscopic_m2_X10000000.png}
\caption{Finite logarithmic mesoscopic check for $m=2$, $X=10^7$. The dashed
curve is the envelope supplied by the unconditional signed majority bound.}
\label{fig:logmajority}
\end{figure}

\section{Conclusions}
\label{sec:conclusions}

Reduced rational dilations decompose into one proportional residue class and
finitely many non-proportional affine classes. Under fixed affine Chowla this
produces the natural microscopic profile
$\lambda(pq)\min(1/q,1/(pc))$ on the positive side and zero on the negative
side. Logarithmic weighting behaves differently: the fixed $X^{-1}$ layer
collapses to a one-sided step, while the elementary logarithmic majority mechanism has threshold
exponent $1/2$ and is effective for every
$c_X=X^{\theta+o(1)}$ with $\theta<1/2$. These unconditional majority
statements distinguish integer from non-integer rational resonances without
using Chowla, while the conditional and unconditional parts of the argument
remain separated explicitly throughout.


With the ordinary two-point theorem of \cite{OpenAI2026} the natural profile
is unconditional, with the rate $(\log X)^{-\kappa}$, because its remainder is
a finite set of fixed binary correlations. At level $k\ge3$ the exact
decomposition reduces every block to a $k'$-point correlation with
$k'\equiv k\pmod2$; blocks with $k'\le2$ are unconditional, and the
left-hand identity shows that blocks with $k'\ge3$ are genuine, for every
large $X$: level three at $(2,3)$ is the ordinary sum
$\sum_{M\equiv3\,(6)}\lambda(M)\lambda(M+1)\lambda(M+3)$. In the bias method
the centering step is controlled by local $U^k$ norms, and cannot be
controlled by the averaged short exponential sums used for two points; the
graph step requires a $k$-linear divisor-graph estimate. Both inputs are
isolated as hypotheses.

\section{Data and reproducibility}
\label{sec:data}

The numerical calculations use only
\path{liouville_int8_N21000000.bin}, containing the first $21{,}000{,}000$
Liouville values as signed int8 with byte $n-1$ equal to $\lambda(n)$. The
binary is deterministically reproducible from
$\lambda(n)=(-1)^{\Omega(n)}$ using the supplied generator. The repository
contains every Python script used for the reported checks, all corresponding
stdout captures, CSV and text outputs, the three manuscript figures as PNG
files, and SHA-256 inventories for the submission files and the complete
repository.

\section{Use of AI tools}
\label{sec:ai}

AI tools were used during drafting, code development, and consistency
checking. All mathematical claims retained in the manuscript were reviewed by
the author. All reported numerical values were regenerated from the supplied
Liouville binary by the supplied Python scripts and checked against the
identities stated in the manuscript. The author verified the resulting
computations and remains responsible for the manuscript and its numerical
content.

\begin{thebibliography}{99}
\raggedright

\bibitem{Abercrombie1995}
A.~G. Abercrombie,
\textit{Beatty sequences and multiplicative number theory},
Acta Arith. 70 (1995), 195--207; DOI: \href{https://doi.org/10.4064/aa-70-3-195-207}{10.4064/aa-70-3-195-207}.

\bibitem{AbercrombieBanksShparlinski2009}
A.~G. Abercrombie, W.~D. Banks, I.~E. Shparlinski,
\textit{Arithmetic functions on Beatty sequences},
Acta Arith. 136 (2009), 81--89; DOI: \href{https://doi.org/10.4064/aa136-1-6}{10.4064/aa136-1-6}.

\bibitem{Frantzikinakis2020}
N. Frantzikinakis,
\textit{Correlations of multiplicative functions along deterministic and
independent sequences},
arXiv:\href{https://arxiv.org/abs/1908.02732}{1908.02732}.

\bibitem{CHRST2024}
D. Crn\v{c}evi\'c, F. Hern\'andez, K. Rizk, K. Sereesuchart, R. Tao,
\textit{On the multiplicative independence between $n$ and
$\lfloor \alpha n\rfloor$},
Acta Arith. 213 (2024), 193--226;
DOI: \href{https://doi.org/10.4064/aa230115-18-2}{10.4064/aa230115-18-2};
arXiv:\href{https://arxiv.org/abs/2211.15830}{2211.15830}.

\bibitem{MRT2015}
K. Matom\"aki, M. Radziwi{\l}{\l}, T. Tao,
\textit{An averaged form of Chowla's conjecture},
Algebra Number Theory 9 (2015), 2167--2196; DOI: \href{https://doi.org/10.2140/ant.2015.9.2167}{10.2140/ant.2015.9.2167}.

\bibitem{Tao2016}
T. Tao,
\textit{The logarithmically averaged Chowla and Elliott conjectures for
two-point correlations},
Forum Math. Pi 4 (2016), e8; DOI: \href{https://doi.org/10.1017/fmp.2016.6}{10.1017/fmp.2016.6};
arXiv:\href{https://arxiv.org/abs/1509.05422}{1509.05422}.

\bibitem{TeravainenWalker2025}
J. Ter\"av\"ainen, A. Walker,
\textit{On a Bohr set analogue of Chowla's conjecture},
Math. Z. 310 (2025), Article 75; DOI: \href{https://doi.org/10.1007/s00209-025-03770-2}{10.1007/s00209-025-03770-2};
arXiv:\href{https://arxiv.org/abs/2303.12574}{2303.12574}.

\bibitem{CasanaPadic}
V. Casa\~na-Giner,
\textit{Exact $p$-adic martingale shells and finite-conductor encodings},
manuscript (padic\_shells), 2026; \textit{Theorems and Their Proofs}, compendium of 25 October 2026, entries 48.0.32--48.0.36, pp.~884--887.

\bibitem{Gowers2001}
W.~T. Gowers,
\textit{A new proof of Szemer\'edi's theorem},
Geom. Funct. Anal. 11 (2001), 465--588;
DOI: \href{https://doi.org/10.1007/s00039-001-0332-9}{10.1007/s00039-001-0332-9}.

\bibitem{MRT2020}
K. Matom\"aki, M. Radziwi{\l}{\l}, T. Tao,
\textit{Fourier uniformity of bounded multiplicative functions in short intervals on average},
Invent. Math. 220 (2020), 1--58;
DOI: \href{https://doi.org/10.1007/s00222-019-00926-w}{10.1007/s00222-019-00926-w};
arXiv:\href{https://arxiv.org/abs/1812.01224}{1812.01224}.

\bibitem{OpenAI2026}
OpenAI,
\textit{Ordinary two-point correlations of multiplicative functions},
preprint, 24 September 2026, 87 pp.; no DOI.

\bibitem{Tao2017}
T. Tao,
\textit{Equivalence of the logarithmically averaged Chowla and Sarnak conjectures},
in: Number Theory -- Diophantine Problems, Uniform Distribution and Applications,
Springer, Cham, 2017, pp.~391--421;
DOI: \href{https://doi.org/10.1007/978-3-319-55357-3_21}{10.1007/978-3-319-55357-3\_21};
arXiv:\href{https://arxiv.org/abs/1605.04628}{1605.04628}.

\bibitem{TaoTeravainen2018}
T. Tao, J. Ter\"av\"ainen,
\textit{Odd order cases of the logarithmically averaged Chowla conjecture},
J. Th\'eor. Nombres Bordeaux 30 (2018), 997--1015;
DOI: \href{https://doi.org/10.5802/jtnb.1062}{10.5802/jtnb.1062};
arXiv:\href{https://arxiv.org/abs/1710.02112}{1710.02112}.

\end{thebibliography}

\end{document}
